\documentclass[12pt,a4paper,UTF8]{ctexart}
\usepackage[a4paper,top=22mm,bottom=22mm,left=22mm,right=22mm,headheight=15pt]{geometry}
\usepackage{amsmath,amssymb,mathtools,bm}
\usepackage{graphicx}
\usepackage{xcolor}
\usepackage{booktabs,longtable,array}
\usepackage{enumitem}
\usepackage{fancyhdr}
\usepackage{caption}
\usepackage{float}
\usepackage{microtype}
\usepackage{setspace}
\usepackage[hidelinks,unicode]{hyperref}
\usepackage[most]{tcolorbox}

\definecolor{ContestBlue}{HTML}{17365D}
\definecolor{SoftBlue}{HTML}{EEF4FA}
\definecolor{SoftGray}{HTML}{F5F6F7}
\hypersetup{pdftitle={第102套},pdfauthor={物理竞赛教学资料}}

\setstretch{1.24}
\setlength{\parindent}{2em}
\setlength{\parskip}{0.35em}
\setlist{itemsep=0.35em,topsep=0.4em}
\captionsetup{font=small,labelfont=bf,justification=centering}
\setcounter{secnumdepth}{0}

\pagestyle{fancy}
\fancyhf{}
\fancyhead[L]{\small 第102套}
\fancyfoot[C]{\thepage}
\renewcommand{\headrulewidth}{0.4pt}

\ctexset{
  section={format=\Large\bfseries\color{ContestBlue},beforeskip=1.2em,afterskip=0.7em},
  subsection={format=\large\bfseries\color{ContestBlue},beforeskip=1.0em,afterskip=0.5em},
  subsubsection={format=\normalsize\bfseries,beforeskip=0.8em,afterskip=0.35em}
}

\makeatletter
\newcommand{\tightlist}{\setlength{\itemsep}{0pt}\setlength{\parskip}{0pt}}
\newcommand{\maxwidth}{\ifdim\Gin@nat@width>\linewidth\linewidth\else\Gin@nat@width\fi}
\newcommand{\maxheight}{\ifdim\Gin@nat@height>0.72\textheight 0.72\textheight\else\Gin@nat@height\fi}
\newcommand{\pandocbounded}[1]{\begingroup\setkeys{Gin}{width=\maxwidth,height=\maxheight,keepaspectratio}#1\endgroup}
\makeatother

\begin{document}

\begin{titlepage}
  \centering
  \vspace*{58mm}
  {\Huge\bfseries\color{ContestBlue} 第102套\par}
  \vfill
  \vspace*{28mm}
\end{titlepage}


\section{第 1 题}\label{ux7b2c-1-ux9898}

白球为均匀球体，质量为 \(m\)。击球后，白球以初速度 \(v_0\)
撞向一个初始静止、质量也为
\(m\)、半径相同的黑球。由于瞄准位置不固定，两球碰撞时的碰撞参数可以任意变化。不计球与球、球与桌面之间的一切摩擦。

\begin{enumerate}
\def\labelenumi{\arabic{enumi}.}
\tightlist
\item
  若碰撞不是完全弹性的，恢复系数为
  \(e=0.5\)，求碰撞后白球可能达到的最大散射角
  \(\theta_{1\max}\)。散射角定义为白球末速度与初速度的夹角。
\item
  将白球换成半径相同、质量为 \(5m\) 的新白球。新白球仍以初速度 \(v_0\)
  撞向初始静止、质量为 \(m\)
  的黑球，且碰撞完全弹性。求碰撞后新白球可能达到的最大散射角
  \(\theta_{2\max}\)。
\end{enumerate}

\subsection{一、题型定位}\label{ux4e00ux9898ux578bux5b9aux4f4d}

斜碰先按法线和切线分解速度，用动量守恒与恢复系数求碰后速度。散射角极值可在质心系速度圆上作切线求出。

\subsection{二、知识点}\label{ux4e8cux77e5ux8bc6ux70b9}

\begin{itemize}
\tightlist
\item
  光滑球斜碰时，冲量只沿两球球心连线方向，切向速度分量不变。
\item
  恢复系数约束的是接触法线方向的分离相对速度与接近相对速度。
\item
  一维碰撞公式与平面速度合成。
\item
  质心系中的完全弹性碰撞、速度圆和切线极值。
\item
  利用均值不等式或求导求散射角极值。
\end{itemize}

\subsection{三、物理图像}\label{ux4e09ux7269ux7406ux56feux50cf}

碰撞发生的一瞬间，两球之间只有沿球心连线的法向冲量。因而入射白球的速度可分成``法向分量''和``切向分量''：前者参与一维碰撞，后者原封不动地保留下来。散射角的变化本质上来自这两个分量在碰后的重新组合。

第二问在实验室系中看似仍需讨论碰撞方向；转到质心系后，重白球的速度大小不变而方向任意，变回实验室系便得到一个固定圆。最大散射角对应从原点向该圆作切线。

\subsection{四、解题入口}\label{ux56dbux89e3ux9898ux5165ux53e3}

设碰撞时球心连线与白球初速度的夹角为 \(\alpha\)，其中
\(0\leq\alpha<\pi/2\)。以球心连线为 \(x'\) 轴、其垂线为 \(y'\)
轴。先只处理 \(x'\) 方向的一维碰撞，再把不变的 \(y'\)
分量加回去。不要一开始就在实验室直角坐标系中列两个方向的动量方程，否则恢复系数所约束的方向容易写错。

\subsection{五、解析}\label{ux4e94ux89e3ux6790}

图 1-1 左侧把入射速度分解到球心连线法向 \(n\) 与切向
\(t\)；右侧画出质量为 \(5m\)
的白球在完全弹性碰撞后速度端点圆及从速度原点所作切线。第（1）问沿左图只更新法向分量，第（2）问沿右图由切线条件求最大偏转角。

\begin{figure}[H]
\centering
\includegraphics[width=0.9\linewidth,height=\textheight,keepaspectratio,alt={图 1-1　斜碰的法向---切向分解，以及重白球在实验室系中的速度圆。}]{../figures/solution-01.pdf}
\caption*{图
1-1　斜碰的法向---切向分解，以及重白球在实验室系中的速度圆。}
\end{figure}

\textbf{（1）等质量、恢复系数 \(e=0.5\)}

按图 1-1 左侧，光滑球碰撞冲量只沿球心连线，所以切向分量保持
\(v_0\sin\alpha\)
不变；恢复系数只约束法向相对速度。先做一维法向碰撞，再把切向分量加回实验室速度。

碰前白球在法向和切向的速度分量分别为

\[v_{0n}=v_0\cos\alpha,\qquad v_{0t}=v_0\sin\alpha.\]

碰撞冲量沿法向，因此白球碰后的切向分量不变：

\[v_t=v_0\sin\alpha. \tag{1}\]

设白球和黑球碰后的法向速度分别为
\(v_n\)、\(u_n\)，取入射法向为正方向。法向动量守恒给出

\[mv_0\cos\alpha=mv_n+mu_n. \tag{2}\]

恢复系数的定义为

\[u_n-v_n=e(v_0\cos\alpha-0). \tag{3}\]

将式（2）除以 \(m\)，与式（3）联立：

\[\begin{aligned}
v_n+u_n&=v_0\cos\alpha,\\
-v_n+u_n&=e v_0\cos\alpha.
\end{aligned}\]

两式相加、相减，分别得到

\[u_n=\frac{1+e}{2}v_0\cos\alpha,\qquad
v_n=\frac{1-e}{2}v_0\cos\alpha. \tag{4}\]

设碰后白球速度与法线的夹角为 \(\alpha'\)。由式（1）、（4），

\[\tan\alpha'=\frac{v_t}{v_n}
=\frac{2\tan\alpha}{1-e}. \tag{5}\]

白球相对初速度的散射角为

\[\theta=\alpha'-\alpha.\]

使用正切差角公式：

\[\begin{aligned}
\tan\theta
&=\frac{\tan\alpha'-\tan\alpha}
{1+\tan\alpha'\tan\alpha}\\
&=\frac{\dfrac{2}{1-e}\tan\alpha-\tan\alpha}
{1+\dfrac{2}{1-e}\tan^2\alpha}\\
&=\frac{(1+e)\tan\alpha}
{1-e+2\tan^2\alpha}.
\end{aligned} \tag{6}\]

代入 \(e=0.5\)，并令 \(x=\tan\alpha>0\)：

\[\tan\theta=\frac{3}{x^{-1}+4x}. \tag{7}\]

除均值不等式外，也可用导数完整检查这是全局最大值。令

\[D(x)=\frac1x+4x,\]

则

\[D'(x)=-\frac1{x^2}+4.\]

驻点条件 \(D'(x)=0\) 给出

\[x=\frac12.\]

而

\[D''(x)=\frac{2}{x^3}>0,\]

故该驻点是 \(D(x)\) 的严格极小值。再检查端点：

\[\lim_{x\to0^+}D(x)=+\infty,\]

\[\lim_{x\to+\infty}D(x)=+\infty.\]

所以 \(x=1/2\) 不只是局部极值，而是 \(x>0\) 上的全局极小点，相应地
\(\theta\) 为全局最大值。

要使 \(\theta\) 最大，只需使分母 \(x^{-1}+4x\) 最小。由均值不等式

\[\frac1x+4x\geq2\sqrt{\frac1x\cdot4x}=4,\]

等号在 \(1/x=4x\)，即 \(x=1/2\) 时成立。因此

\[\tan\theta_{1\max}=\frac34,\]

从而

\[\boxed{\theta_{1\max}=\arctan\frac34\approx36.87^\circ}. \tag{8}\]

这里 \(\alpha=\arctan(1/2)\)
确实位于允许区间内，且碰后法向分量为正，因此极值可实现。

\textbf{（2）质量为 \(5m\) 的白球与质量为 \(m\) 的黑球完全弹性碰撞}

图 1-1
右侧中，白球实验室速度等于质心平动速度与质心系速度的矢量和，故端点落在固定圆上。最大偏转方向是从速度原点到该圆的切线，切点半径与出射速度垂直。

质心速度为

\[\boldsymbol V_C=\frac{5m\boldsymbol v_0+m\boldsymbol 0}{6m}
=\frac56\boldsymbol v_0. \tag{9}\]

碰前重白球相对质心的速度为

\[\boldsymbol v_0'=\boldsymbol v_0-\boldsymbol V_C
=\frac16\boldsymbol v_0.\]

完全弹性碰撞在质心系中保持每个球的速率不变，所以碰后重白球的质心系速度
\(\boldsymbol v'\) 满足

\[|\boldsymbol v'|=\frac16v_0, \tag{10}\]

而其方向可以随碰撞参数改变。回到实验室系，

\[\boldsymbol v=\boldsymbol V_C+\boldsymbol v'.\]

因此 \(\boldsymbol v\) 的末端落在一个圆上：圆心在初速度方向上，距原点
\(5v_0/6\)；圆半径为
\(v_0/6\)。从速度空间原点向该圆作切线，切线与初速度方向的夹角就是最大散射角。由直角三角形

\[\sin\theta_{2\max}
=\frac{v_0/6}{5v_0/6}=\frac15.\]

故

\[\boxed{\theta_{2\max}=\arcsin\frac15\approx11.54^\circ}. \tag{11}\]

速度圆还可推广到任意质量比，以检验本题的几何结论。若入射球质量为
\(Mm\)、靶球质量为 \(m\)，则

\[V_C=\frac{M}{M+1}v_0,\]

而入射球在质心系中的速率为

\[v'_0=v_0-V_C=\frac{1}{M+1}v_0.\]

完全弹性碰撞后，实验室系末速度端点仍在圆上，其圆心距与半径分别为

\[d=\frac{M}{M+1}v_0,\]

\[\rho=\frac{1}{M+1}v_0.\]

当 \(M>1\) 时原点在圆外，最大散射角由切线给出：

\[\sin\theta_{\max}=\frac{\rho}{d}=\frac1M.\]

代入 \(M=5\) 立即回到式（11）。当 \(M=1\) 时圆过原点；当 \(M<1\)
时原点落入圆内，入射球甚至可以出现向后的散射。这组三种几何关系解释了质量比为何决定散射角上限。

第一问的法线---切线分解必须建立在``接触面光滑''这一条件上。碰撞冲量沿两球球心连线，故白球切向分量从碰前到碰后保持不变；法向分量则与黑球共同满足一维动量守恒和恢复系数定义。设入射速度与碰撞法线夹角为
\(\alpha\)，先求出碰后白球相对法线的方向角 \(\alpha'\)：

\[\tan\alpha'=\frac{v_t}{v_n}.\]

再由几何关系求真正的散射角：

\[\theta=\alpha'-\alpha.\]

因此不能直接把 \(v_t/v_n\) 当成 \(\tan\theta\)，也不能把 \(\alpha\) 或
\(\alpha'\) 单独当成散射角。极值条件既可对式（6）的
\(\tan\theta(\alpha)\)
求导，也可在速度空间研究末速度端点轨迹；求得候选角后必须确认两球碰前确实在接近，即法向相对速度为正。

第二问采用质心系的原因是完全弹性碰撞在该系中只改变速度方向、不改变各自速率。重白球实验室末速度端点构成圆心为
\(\boldsymbol V_C\)、半径为 \(v_0/6\)
的圆。最大散射角不是圆上纵坐标最大的点，而是从速度原点向圆作切线的切点，因为散射角是末速度矢量与初速度轴的夹角。质量比增大时圆心更远、半径更小，最大偏角随之减小。

\newpage

\section{第 2 题}\label{ux7b2c-2-ux9898}

一个半径为 \(R\)、质量为 \(M\) 的半圆柱放在水平面上。另一个半径为
\(r\)、质量为 \(m\)
的均匀实心小圆柱从半圆柱顶端由静止开始滚下。小圆柱与半圆柱之间的摩擦因数为
\(\mu_1\)，半圆柱与地面之间的摩擦因数为
\(\mu_2\)。小圆柱圆心相对半圆柱圆心转过的角度记为 \(\theta\)。

\begin{enumerate}
\def\labelenumi{\arabic{enumi}.}
\tightlist
\item
  \(\mu_1\)、\(\mu_2\) 均足够大。求小圆柱离开半圆柱时的角度 \(\theta\)。
\item
  \(\mu_2\) 足够大而 \(\mu_1\)
  较小。求小圆柱与半圆柱之间开始打滑时的角度 \(\theta\)。
\item
  \(\mu_1\) 足够大而 \(\mu_2\) 较小。求半圆柱与地面之间开始打滑的角度
  \(\theta\) 所满足的一元方程。
\item
  取 \(M=2m\)、\(\mu_2=0\)。当 \(\mu_1\)
  足够大时，判断小圆柱能否离开半圆柱；若能，求离开角度。当 \(\mu_1\)
  较小时，求小圆柱与半圆柱之间开始打滑的角度 \(\theta\)
  所满足的一元方程。
\end{enumerate}

\begin{figure}[H]
\centering
\includegraphics[width=0.9\linewidth,height=\textheight,keepaspectratio,alt={图 2-1　小圆柱沿半圆柱外表面滚下，图示角度从竖直方向量起。}]{../figures/question-02.pdf}
\caption*{图 2-1　小圆柱沿半圆柱外表面滚下，图示角度从竖直方向量起。}
\end{figure}

\subsection{一、题型定位}\label{ux4e00ux9898ux578bux5b9aux4f4d-1}

小球在半圆柱上滚动，需同时检查球与柱、柱与地面两处接触。前三问柱固定，第四问柱可平动，须增加水平动量方程，并重新写动能。

\subsection{二、知识点}\label{ux4e8cux77e5ux8bc6ux70b9-1}

\begin{itemize}
\tightlist
\item
  均匀实心圆柱的转动惯量 \(I=\frac12mr^2\)。
\item
  无滑动条件 \(r\omega=(R+r)\dot\theta\)。
\item
  沿圆周的切向方程与指向曲率中心的径向方程。
\item
  脱离条件 \(N=0\)，滑动临界条件 \(|f|=\mu N\)。
\item
  可动基座系统的水平动量守恒、能量守恒和非惯性系惯性力。
\item
  隐式方程的根区间与物理先后次序判断。
\end{itemize}

\subsection{三、物理图像}\label{ux4e09ux7269ux7406ux56feux50cf-1}

小圆柱的圆心沿半径 \(L=R+r\)
的圆弧运动。只要接触存在，半圆柱给小圆柱一个沿外法线的支持力
\(N\)；重力指向圆心的分量逐渐减小，而圆周运动所需向心力逐渐增大，二者竞争决定何时脱离。

静摩擦不是预先等于
\(\mu N\)。在未滑动阶段，它由平动和转动方程共同决定；只有达到极限时才取等号。当地面无摩擦时，半圆柱会反向平移，小圆柱的绝对加速度中必须计入基座加速度。

\subsection{四、解题入口}\label{ux56dbux89e3ux9898ux5165ux53e3-1}

记 \(L=R+r\)。先在``半圆柱固定''模型中一次性求出

\[\dot\theta^2,\qquad \ddot\theta,\qquad N,\qquad f,\]

前三问都从这四个量出发。第四问先由水平动量守恒消去半圆柱速度，再写总动能；不要直接沿用前三问的
\(\dot\theta^2\)。

\subsection{五、解析}\label{ux4e94ux89e3ux6790-1}

图 2-2 左侧画出半圆柱固定时小圆柱所受重力、支持力 \(N\) 和切向静摩擦力
\(f\)；右侧画出地面无摩擦时半圆柱向左速度 \(v\)、小圆柱随动系切向速度
\(L\dot\theta\)
及惯性力。第（1）---（3）步使用左图，第（4）步使用右图重新建立动量与能量关系。

\begin{figure}[H]
\centering
\includegraphics[width=0.9\linewidth,height=\textheight,keepaspectratio,alt={图 2-2　小圆柱的切向---径向受力，以及地面无摩擦时基座与小圆柱的速度关系。}]{../figures/solution-02.pdf}
\caption*{图
2-2　小圆柱的切向---径向受力，以及地面无摩擦时基座与小圆柱的速度关系。}
\end{figure}

记

\[L=R+r,\qquad I=\frac12mr^2.\]

\textbf{（1）两处均不打滑，半圆柱固定}

按图 2-2
左侧，重力切向分量驱动质心运动，静摩擦同时提供转动角加速度；径向方程中的
\(N\) 则决定何时离面。三个方程必须配合纯滚动约束一起使用。

无滑动条件为

\[\omega=\frac{L}{r}\dot\theta. \tag{1}\]

从顶端由静止滚到角度 \(\theta\)，重力势能减少

\[\Delta U=mgL(1-\cos\theta).\]

机械能守恒：

\[mgL(1-\cos\theta)
=\frac12mL^2\dot\theta^2
+\frac12I\left(\frac{L}{r}\dot\theta\right)^2. \tag{2}\]

式（2）中的两项动能分别是小圆柱质心平动动能和绕质心转动动能。转动项逐步化简为

\[\begin{aligned}
K_{\rm rot}
&=\frac12I\omega^2\\
&=\frac12\left(\frac12mr^2\right)
\left(\frac{L}{r}\dot\theta\right)^2\\
&=\frac14mL^2\dot\theta^2.
\end{aligned}\]

质心平动项为

\[K_{\rm tr}=\frac12m(L\dot\theta)^2
=\frac12mL^2\dot\theta^2.\]

因此转动动能恰为平动动能的一半，总动能为

\[K=K_{\rm tr}+K_{\rm rot}
=\frac34mL^2\dot\theta^2.\]

这一比例是均匀实心圆柱的特征；若转动惯量改变，后面的离面角也会随之改变。

代入 \(I=\frac12mr^2\)：

\[\begin{aligned}
mgL(1-\cos\theta)
&=\frac12mL^2\dot\theta^2
+\frac14mL^2\dot\theta^2\\
&=\frac34mL^2\dot\theta^2.
\end{aligned}\]

故

\[\dot\theta^2=\frac{4g}{3L}(1-\cos\theta). \tag{3}\]

沿指向半圆柱圆心的径向列牛顿第二定律。重力的向心分量为
\(mg\cos\theta\)，支持力 \(N\) 指向外侧：

\[mg\cos\theta-N=mL\dot\theta^2. \tag{4}\]

脱离瞬间 \(N=0\)。把式（3）代入式（4）：

\[g\cos\theta=\frac{4g}{3}(1-\cos\theta).\]

因此

\[7\cos\theta=4,\]

即

\[\boxed{\theta_1=\arccos\frac47\approx55.15^\circ}. \tag{5}\]

\textbf{（2）小圆柱与半圆柱之间开始打滑}

在打滑之前仍满足纯滚动。对小圆柱沿切向列式：

\[mg\sin\theta-f=mL\ddot\theta. \tag{6}\]

绕小圆柱圆心列转动方程：

\[fr=I\dot\omega
=\frac12mr^2\frac{L}{r}\ddot\theta,\]

从而

\[f=\frac12mL\ddot\theta. \tag{7}\]

将式（7）代入式（6）：

\[mg\sin\theta=\frac32mL\ddot\theta,\]

所以

\[\ddot\theta=\frac{2g}{3L}\sin\theta,\qquad
f=\frac13mg\sin\theta. \tag{8}\]

再由式（3）、（4）求支持力：

\[\begin{aligned}
N
&=mg\cos\theta-mL\dot\theta^2\\
&=mg\cos\theta-\frac43mg(1-\cos\theta)\\
&=\frac{mg}{3}(7\cos\theta-4).
\end{aligned} \tag{9}\]

开始打滑时 \(f=\mu_1N\)，故

\[\boxed{\sin\theta_2=\mu_1(7\cos\theta_2-4)}. \tag{10}\]

临界式还给出摩擦需求随角度的变化：

\[\frac{f}{N}
=\frac{\sin\theta}{7\cos\theta-4}.\]

在 \(0<\theta<\arccos(4/7)\) 内，分母为正。其导数为

\[\frac{\mathrm d}{\mathrm d\theta}
\left(\frac{\sin\theta}{7\cos\theta-4}\right)
=\frac{7-4\cos\theta}{(7\cos\theta-4)^2}>0.\]

所以所需摩擦因数从 \(0\) 单调增大到 \(+\infty\)。任意有限的 \(\mu_1\)
都对应唯一的首次打滑角，而且该角一定早于无滑动模型的离面角。

将其写成 \(\sin\theta_2-7\mu_1\cos\theta_2=-4\mu_1\)，可得等价的显式形式

\[\boxed{
\theta_2=\arcsin\frac{7\mu_1}{\sqrt{49\mu_1^2+1}}
-\arcsin\frac{4\mu_1}{\sqrt{49\mu_1^2+1}}
}. \tag{11}\]

取这一支是因为物理过程从 \(\theta=0\) 开始，需选取区间
\(0<\theta_2<\arccos(4/7)\) 内首先达到摩擦极限的根。

\textbf{（3）半圆柱与地面之间开始打滑}

小圆柱仍纯滚，故式（8）、（9）继续成立。设地面对半圆柱的水平摩擦力大小为
\(f'\)、竖直支持力为 \(N'\)。半圆柱尚未运动时，水平方向平衡给出

\[f'=N\sin\theta-f\cos\theta. \tag{12}\]

竖直方向平衡给出

\[N'=Mg+N\cos\theta+f\sin\theta. \tag{13}\]

把式（8）、（9）代入式（12）：

\[\begin{aligned}
f'
&=\frac{mg}{3}(7\cos\theta-4)\sin\theta
-\frac{mg}{3}\sin\theta\cos\theta\\
&=mg\sin\theta\left(2\cos\theta-\frac43\right).
\end{aligned} \tag{14}\]

同理，

\[\begin{aligned}
N'
&=Mg+\frac{mg}{3}(7\cos\theta-4)\cos\theta
+\frac{mg}{3}\sin^2\theta\\
&=Mg+mg\left(2\cos^2\theta-\frac43\cos\theta+\frac13\right).
\end{aligned} \tag{15}\]

地面开始打滑时 \(|f'|=\mu_2N'\)。用带符号的 \(f'\)
表示地面所需摩擦，其方向在 \(\cos   heta=2/3\)
时改变。两段都须保留绝对值，首次打滑角满足

\[\boxed{
mg\sin\theta\left|2\cos\theta-\frac43\right|
=\mu_2\left[Mg+mg\left(2\cos^2\theta-\frac43\cos\theta+\frac13\right)\right]
}. \tag{16}\]

式（16）中的符号也可以从初始邻域检查。令 \(\theta\ll1\)，则

\[f'\approx mg\theta\left(2-\frac43\right)
=\frac23mg\theta>0,\]

而

\[N'\approx Mg+mg\left(2-\frac43+\frac13\right)
=(M+m)g.\]

因此在运动刚开始时，地面摩擦确实取式（12）规定的正方向。随着角度增大，\(f'\)
可能变号，所以若题目继续追踪更晚过程，必须恢复绝对值；本问只需寻找第一个达到极限的根。

求根时应取 \(0<\theta<\arccos(4/7)\)
内的最小正根；若该区间无根，说明在小圆柱离开前地面不会打滑。

\textbf{（4）\(M=2m\)、地面无摩擦}

转看图 2-2
右侧。地面无水平外力时，半圆柱会反冲，故小圆柱的地面系水平速度是相对基座切向速度的水平分量减去基座速度
\(v\)；不能沿用固定基座时的动能式。

设半圆柱向左的速率为 \(v>0\)。小圆柱圆心相对半圆柱的切向速度为
\(L\dot\theta\)，其水平分量向右为
\(L\dot\theta\cos\theta\)，故小圆柱相对地面的水平速度为

\[v_{mx}=L\dot\theta\cos\theta-v.\]

系统初始水平动量为零，因此

\[Mv=m(L\dot\theta\cos\theta-v). \tag{17}\]

代入 \(M=2m\)：

\[\boxed{v=\frac13L\dot\theta\cos\theta}. \tag{18}\]

总机械能守恒：

\[\begin{aligned}
mgL(1-\cos\theta)
={}&\frac12Mv^2\\
&+\frac12m\left[(L\dot\theta\sin\theta)^2
+(L\dot\theta\cos\theta-v)^2\right]\\
&+\frac12I\left(\frac{L}{r}\dot\theta\right)^2.
\end{aligned} \tag{19}\]

将 \(M=2m\)、\(I=\frac12mr^2\) 和式（18）逐项代入，右端化为

\[\begin{aligned}
K
&=\frac12mL^2\dot\theta^2
\left[\frac{2\cos^2\theta}{9}
+\sin^2\theta+\frac{4\cos^2\theta}{9}+\frac12\right]\\
&=\frac{mL^2\dot\theta^2}{12}(9-2\cos^2\theta)\\
&=\frac{mL^2\dot\theta^2}{12}(8-\cos2\theta).
\end{aligned}\]

因此

\[\boxed{
\dot\theta^2
=\frac{24g\sin^2(\theta/2)}{L(8-\cos2\theta)}
=\frac{12g(1-\cos\theta)}{L(9-2\cos^2\theta)}
}. \tag{20}\]

从第一种写法到第二种写法使用了

\[\sin^2\frac\theta2=\frac{1-\cos\theta}{2},\]

以及

\[8-\cos2\theta
=8-(2\cos^2\theta-1)
=9-2\cos^2\theta.\]

这一步把所有三角函数统一成
\(c=\cos\theta\)，便于后面把离面条件化为代数方程。

对式（20）关于时间求导。使用
\(2\dot\theta\ddot\theta=(\mathrm d\dot\theta^2/\mathrm d\theta)\dot\theta\)，约去非零的
\(\dot\theta\)：

\[\boxed{
\ddot\theta
=\frac{6g\sin\theta(9+2\cos^2\theta-4\cos\theta)}
{L(9-2\cos^2\theta)^2}
}. \tag{21}\]

为看清式（21）的来源，记

\[F(\theta)=\frac{12g}{L}
\frac{1-\cos\theta}{9-2\cos^2\theta}.\]

则

\[2\ddot\theta=F'(\theta).\]

用商法则：

\[F'(\theta)=\frac{12g}{L}
\frac{\sin\theta(9-2\cos^2\theta)
-4\cos\theta\sin\theta(1-\cos\theta)}
{(9-2\cos^2\theta)^2}.\]

提取 \(\sin\theta\) 并合并括号：

\[9-2\cos^2\theta-4\cos\theta+4\cos^2\theta
=9+2\cos^2\theta-4\cos\theta.\]

再除以 \(2\)，即得到式（21）。

由式（18），半圆柱向左的加速度为

\[\begin{aligned}
a
&=\dot v
=\frac{L}{3}\left(\cos\theta\,\ddot\theta
-\sin\theta\,\dot\theta^2\right)\\
&=-\frac{2g\sin\theta(2\cos^3\theta-27\cos\theta+18)}
{(9-2\cos^2\theta)^2}.
\end{aligned} \tag{22}\]

直接对式（18）求导，可以独立核验式（22）的每个系数和符号。

在随半圆柱向左加速的参考系中，小圆柱受到水平向右的惯性力
\(ma\)。沿向心方向列式：

\[mg\cos\theta-N-ma\sin\theta=mL\dot\theta^2. \tag{23}\]

代入式（20）、（22）并整理，得到

\[\frac{N}{mg}
=\frac{-2c^3-12c^2+135c-72}{(9-2c^2)^2},
\qquad c=\cos\theta. \tag{24}\]

脱离条件 \(N=0\)，即

\[2c^3+12c^2-135c+72=0. \tag{25}\]

在 \(0<c<1\) 内唯一物理解为

\[c=0.5643008285\ldots,\]

所以小圆柱能够离开半圆柱，离开角为

\[\boxed{\theta_{\rm off}=\arccos(0.5643008285\ldots)
\approx55.64^\circ}. \tag{26}\]

再求两者之间维持纯滚所需的摩擦力。沿切向列式：

\[mg\sin\theta+ma\cos\theta-f=mL\ddot\theta. \tag{27}\]

代入式（20）---（22）可得

\[\frac{f}{mg}
=\frac{3\sin\theta(2c^2-4c+9)}{(9-2c^2)^2}. \tag{28}\]

开始打滑时 \(f=\mu_1N\)。由式（24）、（28），临界角满足

\[\boxed{
3\sin\theta(2c^2-4c+9)
=\mu_1(-2c^3-12c^2+135c-72),
\qquad c=\cos\theta
}. \tag{29}\]

应取 \(0<\theta<\theta_{\rm off}\)
内最先出现的正根。若无根，则在离开之前始终不打滑。

前三问中半圆柱相对地面静止，小圆柱纯滚时有 \(v_C=(R+r)\dot\theta\)
与自转角速度之间的无滑动约束。机械能同时包含质心平动和绕自身轴转动；对均匀实心圆柱
\(I=mr^2/2\)，所以转动动能不能遗漏。由能量先求
\(\dot\theta^2(\theta)\)，再对时间求导得到
\(\ddot\theta\)，随后才能分别进入径向和切向受力方程。

求摩擦时要先规定正方向。摩擦力的实际方向不是永远``沿斜面向上''，而由维持纯滚所需的平动---转动协调决定；算出
\(f\) 为负只表示实际方向与假设相反。判断打滑用
\(|f|=\mu_1N\)，判断离面用
\(N=0\)，两个条件要比较谁先达到。地面摩擦条件同样要先求半圆柱受到的小圆柱作用力，再由水平平衡得到地面所需摩擦，并与
\(\mu_2N_g\) 比较。

第四问中基座可以水平运动，选择广义坐标 \(X\) 和 \(\theta\)
后，水平动量守恒先消去 \(\dot X\)，能量方程再给出
\(\dot\theta\)。对时间求导时 \(X\)
的加速度会进入小圆柱的径向、切向方程，因此固定基座所得的 \(N\)、\(f\)
不能直接套用。最终的一元方程只给出候选角，仍应在
\(0<\theta<\theta_{\rm off}\) 中取首次根，并核验在此之前 \(N>0\)。

\newpage

\section{第 3 题}\label{ux7b2c-3-ux9898}

在 \(xOy\) 平面内有一半径为 \(r\)、圆心位于原点 \(O\)
的细圆环。圆环上一点的位置由其半径与 \(x\) 轴正方向的夹角 \(\theta\)
表示，其中 \(0\leq\theta\leq2\pi\)。

\begin{enumerate}
\def\labelenumi{\arabic{enumi}.}
\tightlist
\item
  圆环均匀带电，线电荷密度为常量 \(\lambda\)。

  \begin{enumerate}
  \def\labelenumii{\arabic{enumii}.}
  \tightlist
  \item
    求原点处的电势 \(\varphi_0\)。
  \item
    求原点附近三维空间中的电场，并保留到坐标小量的二阶精度。
  \end{enumerate}
\item
  圆环的线电荷密度为
  \[\lambda(\theta)=\lambda_0\frac{\theta}{2\pi},\qquad 0\leq\theta\leq2\pi.\]

  \begin{enumerate}
  \def\labelenumii{\arabic{enumii}.}
  \tightlist
  \item
    求原点处的电势 \(\varphi_1\) 和电场 \(\boldsymbol E_0\)。
  \item
    求原点附近三维空间中的电场，并保留到坐标小量的一阶。
  \end{enumerate}
\end{enumerate}

\subsection{一、题型定位}\label{ux4e00ux9898ux578bux5b9aux4f4d-2}

求带电圆环附近的电场。均匀带电部分可用对称性；线密度随角度变化时，展开库仑核，再逐项做角积分。

\subsection{二、知识点}\label{ux4e8cux77e5ux8bc6ux70b9-2}

\begin{itemize}
\tightlist
\item
  线电荷的电势积分与电场 \(\boldsymbol E=-\nabla\varphi\)。
\item
  二项式展开 \((1-\delta)^{-1/2}\)。
\item
  先展开后积分的适用条件 \(\rho=\sqrt{x^2+y^2+z^2}\ll r\)。
\item
  三角函数的周期积分及含 \(\theta\) 权重的分部积分。
\item
  对称性对展开阶数的限制。
\end{itemize}

\subsection{三、物理图像}\label{ux4e09ux7269ux7406ux56feux50cf-2}

场点靠近圆心时，圆环上每个电荷元到场点的距离都接近
\(r\)。均匀圆环没有偶极矩，因此原点场强为零，最低非平凡修正是四极型；非均匀圆环则已有沿
\(+y\) 方向的原点电场，并在邻域中叠加线性梯度。

\subsection{四、解题入口}\label{ux56dbux89e3ux9898ux5165ux53e3-2}

对圆环元 \(\boldsymbol r'=r(\cos\theta,\sin\theta,0)\) 和场点
\(\boldsymbol x=(x,y,z)\)，统一写出

\[|\boldsymbol r'-\boldsymbol x|
=r\left[1-2\frac{x\cos\theta+y\sin\theta}{r}
+\frac{x^2+y^2+z^2}{r^2}\right]^{1/2}.\]

随后只需展开其倒数。先明确需要的电场阶数，再决定电势多保留一阶：求一阶电场，电势必须保留到二阶。

\subsection{五、解析}\label{ux4e94ux89e3ux6790-2}

图 3-1 左图画出环面内的源点 \(Q\) 和场点投影
\(P_0(x,y,0)\)，右图标出离开环面的高度 \(z\)。源点矢量为
\(\boldsymbol r'=r(\cos\theta,\sin\theta,0)\)，场点矢量为
\(\boldsymbol\rho=(x,y,z)\)。后面的
\(A=\hat{\boldsymbol r}'\cdot\boldsymbol\rho=x\cos\theta+y\sin\theta\)
是标量投影，因而
\(\boldsymbol r'\cdot\boldsymbol\rho=rA\)，\(\rho^2=x^2+y^2+z^2\)。

\begin{figure}[H]
\centering
\includegraphics[width=0.788\linewidth,height=\textheight,keepaspectratio,alt={图 3-1　圆环电荷元、场点位移与近场展开中的投影量。}]{../figures/solution-03.pdf}
\caption*{图 3-1　圆环电荷元、场点位移与近场展开中的投影量。}
\end{figure}

记

\[A=x\cos\theta+y\sin\theta,\qquad
\rho^2=x^2+y^2+z^2.\]

则

\[\frac1{|\boldsymbol r'-\boldsymbol x|}
=\frac1r\left(1-\frac{2A}{r}+\frac{\rho^2}{r^2}\right)^{-1/2}. \tag{1}\]

利用

\[ (1-u)^{-1/2}=1+\frac12u+\frac38u^2+O(u^3),\]

并只保留坐标的二次项。令 \(u=2A/r-\rho^2/r^2\)，注意 \(A\)
是一阶量、\(\rho^2\) 是二阶量，则

\[\boxed{
\frac1{|\boldsymbol r'-\boldsymbol x|}
=\frac1r\left[1+\frac{A}{r}
-\frac{\rho^2}{2r^2}
+\frac{3A^2}{2r^2}
\right]+O\!\left(\frac{\rho^3}{r^4}\right)
}. \tag{2}\]

式（2）中的阶数筛选可以逐项说明。由

\[u=\frac{2A}{r}-\frac{\rho^2}{r^2},\]

一阶项为

\[\frac12u=\frac{A}{r}-\frac{\rho^2}{2r^2}.\]

平方项为

\[u^2=\frac{4A^2}{r^2}
-\frac{4A\rho^2}{r^3}
+\frac{\rho^4}{r^4}.\]

其中 \(A\rho^2\) 是三阶量，\(\rho^4\)
是四阶量，本题二阶展开都应舍去。因此

\[\frac38u^2=\frac{3A^2}{2r^2}+O(\rho^3/r^3),\]

与一阶项合并后正好得到式（2）。如果不先标注 \(A\) 与 \(\rho^2\)
的阶数，很容易错误保留 \(A\rho^2\)。

\textbf{（1.1）均匀圆环的原点电势}

在图 3-1 的原点 \(P=0\) 时，每个电荷元到场点的距离都等于
\(r\)，所以积分只剩总电荷；从下一小问起才把 \(P\) 视为小量并展开。

电荷元为
\(\mathrm dq=\lambda r\,\mathrm d\theta\)。原点到每个电荷元的距离均为
\(r\)，故

\[\begin{aligned}
\varphi_0
&=\frac1{4\pi\varepsilon_0}
\int_0^{2\pi}\frac{\lambda r\,\mathrm d\theta}{r}\\
&=\frac{\lambda}{4\pi\varepsilon_0}(2\pi).
\end{aligned}\]

所以

\[\boxed{\varphi_0=\frac{\lambda}{2\varepsilon_0}}. \tag{3}\]

\textbf{（1.2）均匀圆环原点附近的电场}

按图 3-1，先把 \(|\boldsymbol r'-\boldsymbol\rho|^{-1}\)
展开到电势二阶，再对角度积分。由于电场是电势的梯度，若只把电势保留到一阶，就无法得到所需的一阶电场。

把式（2）代入电势积分。用到

\[\int_0^{2\pi}\cos\theta\,\mathrm d\theta
=\int_0^{2\pi}\sin\theta\,\mathrm d\theta=0,\]

\[\int_0^{2\pi}\cos^2\theta\,\mathrm d\theta
=\int_0^{2\pi}\sin^2\theta\,\mathrm d\theta=\pi,\]

\[\int_0^{2\pi}\sin\theta\cos\theta\,\mathrm d\theta=0.\]

于是

\[\int_0^{2\pi}A^2\,\mathrm d\theta
=\pi(x^2+y^2).\]

这一积分来自完整展开

\[A^2=x^2\cos^2\theta
+2xy\sin\theta\cos\theta
+y^2\sin^2\theta.\]

逐项积分：

\[\int_0^{2\pi}x^2\cos^2\theta\,\mathrm d\theta=\pi x^2,\]

\[\int_0^{2\pi}2xy\sin\theta\cos\theta\,\mathrm d\theta=0,\]

\[\int_0^{2\pi}y^2\sin^2\theta\,\mathrm d\theta=\pi y^2.\]

交叉项消失正是轴对称性在代数中的表现。

电势为

\[\begin{aligned}
\varphi
&=\frac{\lambda}{4\pi\varepsilon_0}
\int_0^{2\pi}\left[1+\frac{A}{r}
-\frac{\rho^2}{2r^2}+\frac{3A^2}{2r^2}\right]\mathrm d\theta\\
&=\frac{\lambda}{2\varepsilon_0}
+\frac{\lambda}{8\varepsilon_0r^2}(x^2+y^2-2z^2)
+O\!\left(\frac{\lambda}{\varepsilon_0}\frac{\rho^4}{r^4}\right).
\end{aligned} \tag{4}\]

均匀圆环关于原点反演对称，电势只含偶次项；因此电场在一次项之后没有二次项，下一项是三次。取负梯度：

\[\boxed{
\boldsymbol E
=\frac{\lambda}{4\varepsilon_0r^2}
\left(-x\,\hat{\boldsymbol x}-y\,\hat{\boldsymbol y}
+2z\,\hat{\boldsymbol z}\right)+O\!\left(\frac{\lambda\rho^3}{\varepsilon_0r^4}\right)
}. \tag{5}\]

这已经达到``保留到坐标二阶精度''：所有二次电场项因对称性严格为零。

\textbf{（2.1）非均匀圆环的原点电势与电场}

现在

\[\lambda(\theta)=\lambda_0\frac{\theta}{2\pi},\qquad
\mathrm dq=\lambda_0\frac{\theta}{2\pi}r\,\mathrm d\theta.\]

原点电势为

\[\begin{aligned}
\varphi_1
&=\frac1{4\pi\varepsilon_0}
\int_0^{2\pi}\frac{\lambda_0\theta r}{2\pi r}\,\mathrm d\theta\\
&=\frac{\lambda_0}{8\pi^2\varepsilon_0}
\left[\frac{\theta^2}{2}\right]_0^{2\pi}
=\frac{\lambda_0}{4\varepsilon_0}. 
\end{aligned}\]

所以

\[\boxed{\varphi_1=\frac{\lambda_0}{4\varepsilon_0}}. \tag{6}\]

原点处，正电荷元产生的电场方向从电荷元指向原点：

\[\mathrm d\boldsymbol E
=-\frac{\lambda_0}{8\pi^2\varepsilon_0r}
\theta(\cos\theta,\sin\theta,0)\,\mathrm d\theta.\]

分部积分得到

\[\int_0^{2\pi}\theta\cos\theta\,\mathrm d\theta=0,\qquad
\int_0^{2\pi}\theta\sin\theta\,\mathrm d\theta=-2\pi.\]

两个分部积分分别为

\[\begin{aligned}
\int_0^{2\pi}\theta\cos\theta\,\mathrm d\theta
&=\left.\theta\sin\theta\right|_0^{2\pi}
-\int_0^{2\pi}\sin\theta\,\mathrm d\theta\\
&=0,
\end{aligned}\]

以及

\[\begin{aligned}
\int_0^{2\pi}\theta\sin\theta\,\mathrm d\theta
&=\left.-\theta\cos\theta\right|_0^{2\pi}
+\int_0^{2\pi}\cos\theta\,\mathrm d\theta\\
&=-2\pi.
\end{aligned}\]

因此原点电场只有 \(y\) 分量。由于 \(\mathrm d\boldsymbol E\)
前面还有整体负号，积分为 \(-2\pi\) 最终反而给出 \(+y\)
方向，两个负号不能漏掉。

故

\[\boxed{
\boldsymbol E_0=\frac{\lambda_0}{4\pi\varepsilon_0r}
\hat{\boldsymbol y}
}. \tag{7}\]

\textbf{（2.2）非均匀圆环原点附近的一阶电场}

将式（2）乘以带权因子 \(\lambda_0\theta/(2\pi)\)
后积分。除前述积分外，还需

\[\begin{aligned}
\int_0^{2\pi}\theta\,\mathrm d\theta&=2\pi^2,\\
\int_0^{2\pi}\theta\cos^2\theta\,\mathrm d\theta&=\pi^2,\\
\int_0^{2\pi}\theta\sin^2\theta\,\mathrm d\theta&=\pi^2,\\
\int_0^{2\pi}\theta\sin\theta\cos\theta\,\mathrm d\theta&=-\frac\pi2.
\end{aligned}\]

逐项积分后得到二阶电势

\[\boxed{
\varphi(x,y,z)=\frac{\lambda_0}{\varepsilon_0}
\left[
\frac14-\frac{y}{4\pi r}
+\frac{\pi(x^2+y^2-2z^2)-3xy}{16\pi r^2}
\right]+O\!\left(\frac{\lambda_0}{\varepsilon_0}\frac{\rho^3}{r^3}\right)
}. \tag{8}\]

分别对 \(x,y,z\) 求偏导，并加负号：

例如 \(x\) 分量逐步为

\[\frac{\partial\varphi}{\partial x}
=\frac{\lambda_0}{\varepsilon_0}
\frac{2\pi x-3y}{16\pi r^2},\]

所以

\[E_x=-\frac{\partial\varphi}{\partial x}
=\frac{\lambda_0}{\varepsilon_0}
\frac{3y-2\pi x}{16\pi r^2}.\]

同理，

\[\frac{\partial\varphi}{\partial y}
=\frac{\lambda_0}{\varepsilon_0}
\left[-\frac1{4\pi r}
+\frac{2\pi y-3x}{16\pi r^2}\right],\]

\[\frac{\partial\varphi}{\partial z}
=-\frac{\lambda_0}{\varepsilon_0}\frac{z}{4r^2}.\]

逐项取负梯度后，得到

\[\boxed{
\begin{aligned}
E_x&=\frac{\lambda_0}{\varepsilon_0}
\frac{3y-2\pi x}{16\pi r^2},\\
E_y&=\frac{\lambda_0}{\varepsilon_0}
\left(\frac1{4\pi r}+\frac{3x-2\pi y}{16\pi r^2}\right),\\
E_z&=\frac{\lambda_0}{\varepsilon_0}\frac{z}{4r^2}.
\end{aligned}
} \tag{9}\]

令 \(x=y=z=0\)，式（9）退化为式（7），这是常量项检查。

把场点写为 \(\boldsymbol\rho=(x,y,z)\)、源点写为
\(\boldsymbol r'=r(\cos\theta,\sin\theta,0)\) 后，库仑核为

\[\begin{aligned}
\frac1{|\boldsymbol r'-\boldsymbol\rho|}
&=\frac1r\left(1+u\right)^{-1/2},\\
u&=-2\frac{\boldsymbol r'\cdot\boldsymbol\rho}{r^2}
+\frac{\rho^2}{r^2}.
\end{aligned}\]

令括号中的小量为 \(u\)，使用
\((1+u)^{-1/2}=1-u/2+3u^2/8+\cdots\)。若要求电场保留到坐标的一阶，电势必须展开到二阶，因为求梯度会降低一次阶数；这正是本题最容易出现的``少展开一阶''错误。均匀环中所有一阶角积分因对称性为零，二阶项只能组合成
\(x^2+y^2-2z^2\)，其迹为零也与无电荷区域满足拉普拉斯方程一致。

非均匀权重 \(\theta/(2\pi)\)
不是周期光滑函数，不能直接用均匀环的奇偶性删项。应把

\[\int\theta\cos\theta\,\mathrm d\theta,
\qquad
\int\theta\sin\theta\,\mathrm d\theta\]

等积分逐项算出，并保留端点项。得到电势后再统一求负梯度，比直接从三分量电场积分更不容易漏项。最后检查原点电势具有伏特量纲，原点电场方向与电荷较密集的一侧相反。

\newpage

\section{第 4 题}\label{ux7b2c-4-ux9898}

一个半径为 \(R\)、均匀带电量 \(+Q_0\)（\(Q_0>0\)）的绝缘球壳，以角速度

\[\boldsymbol\omega=-\omega\,\hat{\boldsymbol x}\]

匀速转动。

\begin{enumerate}
\def\labelenumi{\arabic{enumi}.}
\tightlist
\item
  电流回路的磁矩定义为
  \(\mathrm d\boldsymbol m=I\,\mathrm d\boldsymbol S\)，其中 \(I\)
  为电流，\(\mathrm d\boldsymbol S\)
  的大小为回路面积、方向由右手定则确定。求转动球壳的总磁矩
  \(\boldsymbol m_0\)。
\item
  写出电偶极子 \(\boldsymbol p\) 在球坐标 \((r,\theta)\)
  处所产生电场的径向分量 \(E_r\) 与极角分量 \(E_\theta\)。
\item
  利用对应关系
  \(1/(4\pi\varepsilon_0)\leftrightarrow\mu_0/(4\pi)\)、\(\boldsymbol p\leftrightarrow\boldsymbol m\)，求远区
  \(r\gg R\) 处的磁感应强度 \(\boldsymbol B\)。角度 \(\theta\) 从 \(+x\)
  轴量起。
\item
  在距球心 \(r\gg R\) 处放置电量为 \(-q\)（\(q>0\)）、质量为 \(m\)
  的点电荷。若它在 \(x=\text{常量}\) 的平面内绕 \(x\)
  轴做匀速圆周运动，求轨道极角 \(\theta\) 及角速度
  \(\omega'\)。忽略辐射、阻力和重力。
\end{enumerate}

\begin{figure}[H]
\centering
\includegraphics[width=0.9\linewidth,height=\textheight,keepaspectratio,alt={图 4-1　带电球壳绕负 x 方向转动，负点电荷在远区绕 x 轴做圆周运动。}]{../figures/question-04.pdf}
\caption*{图 4-1　带电球壳绕负 x 方向转动，负点电荷在远区绕 x
轴做圆周运动。}
\end{figure}

\subsection{一、题型定位}\label{ux4e00ux9898ux578bux5b9aux4f4d-3}

先将旋转球壳拆成纬圈求磁矩，再计算磁偶极远场。粒子的定常圆轨道要分别满足轴向和绕轴径向的受力条件，并检查速度是否属于非相对论范围。

\subsection{二、知识点}\label{ux4e8cux77e5ux8bc6ux70b9-3}

\begin{itemize}
\tightlist
\item
  面电荷密度与旋转纬圈电流。
\item
  电流环磁矩 \(m=IS\) 及矢量方向。
\item
  电偶极与磁偶极的远场公式。
\item
  球坐标基矢到直角坐标分量的转换。
\item
  洛伦兹力
  \(\boldsymbol F=q(\boldsymbol E+\boldsymbol v\times\boldsymbol B)\)。
\item
  绕固定轴圆周运动的轴向平衡与径向向心条件。
\end{itemize}

\subsection{三、物理图像}\label{ux4e09ux7269ux7406ux56feux50cf-3}

把球壳切成无数个垂直于 \(x\)
轴的窄纬圈。每个纬圈都是一个电流环，磁矩都沿角速度方向，积分后得到总磁矩。远处只需保留磁偶极项。

负电荷绕 \(x\) 轴运动时没有 \(x\)
方向加速度，因此电场轴向分量必须由磁力轴向分量抵消；垂直于 \(x\)
轴的合力则提供向心力。两个分量方程共同确定 \(\omega'\) 与 \(\theta\)。

\subsection{四、解题入口}\label{ux56dbux89e3ux9898ux5165ux53e3-3}

以 \(+x\) 轴为极轴定义 \(\theta\)。先暂时令磁矩的代数分量为 \(m_d\)，即
\(\boldsymbol m_0=m_d\hat{\boldsymbol x}\)；本题
\(m_d<0\)。这样偶极场公式的结构保持不变，方向信息由 \(m_d\)
的负号统一承担。

\subsection{五、解析}\label{ux4e94ux89e3ux6790-3}

图 4-2 左侧把带电球壳分解成绕 \(x\)
轴的纬圈电流并叠加磁矩；右侧在负电荷的纬圈圆轨道上分解电场力、磁洛伦兹力与速度方向。第（1）步使用左图求总磁矩，第（2）、（3）步建立远区偶极场，第（4）步按右图令轴向合力为零、横向合力提供向心力。

\begin{figure}[H]
\centering
\includegraphics[width=0.851\linewidth,height=\textheight,keepaspectratio,alt={图 4-2　球壳纬圈磁矩积分，以及负电荷圆轨道上的电场、磁场与洛伦兹力分解。}]{../figures/solution-04.pdf}
\caption*{图
4-2　球壳纬圈磁矩积分，以及负电荷圆轨道上的电场、磁场与洛伦兹力分解。}
\end{figure}

\textbf{（1）球壳的磁矩}

按图 4-2 左侧，极角为 \(\theta\) 的窄纬带既有面积电荷，也以角速度
\(\omega\) 形成电流环；其环面积还含
\(\sin^2\theta\)。球壳带正电，所以每个纬圈的磁矩与角速度同向，均沿
\(-x\)。

球壳面电荷密度为

\[\sigma=\frac{Q_0}{4\pi R^2}. \tag{1}\]

取极角在 \(\theta\) 到 \(\theta+\mathrm d\theta\) 的窄纬带。其面积为

\[\mathrm dA=2\pi R^2\sin\theta\,\mathrm d\theta,\]

所带电量为

\[\mathrm dQ=\sigma\,\mathrm dA
=\frac{Q_0}{2}\sin\theta\,\mathrm d\theta. \tag{2}\]

转动周期为 \(2\pi/\omega\)，故该纬圈形成的电流大小为

\[\mathrm dI=\frac{\mathrm dQ}{2\pi/\omega}
=\frac{Q_0\omega}{4\pi}\sin\theta\,\mathrm d\theta. \tag{3}\]

纬圈半径为 \(R\sin\theta\)，面积为 \(\pi R^2\sin^2\theta\)。磁矩大小

\[\mathrm dm=\mathrm dI\,\pi R^2\sin^2\theta
=\frac{Q_0\omega R^2}{4}\sin^3\theta\,\mathrm d\theta. \tag{4}\]

所有纬圈磁矩方向均沿 \(\boldsymbol\omega\)，即 \(-x\) 方向。又

\[\int_0^\pi\sin^3\theta\,\mathrm d\theta=\frac43,\]

该积分可由 \(u=\cos\theta\) 直接得到：

\[\begin{aligned}
\int_0^\pi\sin^3\theta\,\mathrm d\theta
&=\int_0^\pi(1-\cos^2\theta)\sin\theta\,\mathrm d\theta\\
&=\int_{-1}^{1}(1-u^2)\,\mathrm du\\
&=\left.u-\frac{u^3}{3}\right|_{-1}^{1}\\
&=\frac43.
\end{aligned}\]

所以磁矩大小为

\[m_0=\frac{Q_0\omega R^2}{4}\cdot\frac43
=\frac13Q_0\omega R^2.\]

正电荷随球壳沿 \(\boldsymbol\omega\)
方向旋转，常规电流方向与电荷运动方向相同，故磁矩方向与
\(\boldsymbol\omega\) 同向，即沿 \(-x\)。

故

\[\boxed{
\boldsymbol m_0=-\frac13Q_0\omega R^2\hat{\boldsymbol x}
}. \tag{5}\]

\textbf{（2）电偶极子的场}

若电偶极矩沿极轴正方向，则

\[\boxed{
E_r=\frac1{4\pi\varepsilon_0}\frac{2p\cos\theta}{r^3},
\qquad
E_\theta=\frac1{4\pi\varepsilon_0}\frac{p\sin\theta}{r^3}
}. \tag{6}\]

\textbf{（3）磁偶极远场}

令

\[m_d=-\frac13Q_0\omega R^2,\]

则 \(\boldsymbol m_0=m_d\hat{\boldsymbol x}\)。由电---磁偶极对应关系，

\[\boxed{
\boldsymbol B
=\frac{\mu_0m_d}{4\pi r^3}
\left(2\cos\theta\,\hat{\boldsymbol r}
+\sin\theta\,\hat{\boldsymbol\theta}\right)
}. \tag{7}\]

在图示 \(xOy\) 平面内，

\[\hat{\boldsymbol r}=\cos\theta\hat{\boldsymbol x}
+\sin\theta\hat{\boldsymbol y},\]

\[\hat{\boldsymbol\theta}=-\sin\theta\hat{\boldsymbol x}
+\cos\theta\hat{\boldsymbol y}.\]

代入式（7）：

先分别求 \(x\)、\(y\) 分量：

\[\begin{aligned}
B_x
&=\frac{\mu_0m_d}{4\pi r^3}
\left(2\cos\theta\cos\theta
-\sin\theta\sin\theta\right)\\
&=\frac{\mu_0m_d}{4\pi r^3}(3\cos^2\theta-1),
\end{aligned}\]

其中用了 \(\sin^2\theta=1-\cos^2\theta\)。同理，

\[\begin{aligned}
B_y
&=\frac{\mu_0m_d}{4\pi r^3}
\left(2\cos\theta\sin\theta
+\sin\theta\cos\theta\right)\\
&=\frac{3\mu_0m_d}{4\pi r^3}
\sin\theta\cos\theta.
\end{aligned}\]

于是

\[\boxed{
\boldsymbol B=\frac{\mu_0m_d}{4\pi r^3}
\left[(3\cos^2\theta-1)\hat{\boldsymbol x}
+3\sin\theta\cos\theta\hat{\boldsymbol y}\right]
}. \tag{8}\]

\textbf{（4）负电荷的定常圆轨道}

转看图 4-2 右侧。负电荷的电场力与局部电场反向，磁力由
\(-q\bm v\times\bm B\) 给出；沿 \(x\)
轴分量必须相消，垂直轴的合力才可作为半径 \(r\sin\theta\) 的向心力。

设点电荷位于 \(xOy\) 平面中图示位置，绕 \(x\) 轴的轨道半径为

\[\rho=r\sin\theta.\]

取图示环绕方向的角速度大小为 \(\omega'>0\)，则在所选位置

\[\boldsymbol v=-\omega'\rho\,\hat{\boldsymbol z}.\]

球壳在远处的电场仍是点电荷场：

\[\boldsymbol E=\frac{Q_0}{4\pi\varepsilon_0r^2}
(\cos\theta\hat{\boldsymbol x}+\sin\theta\hat{\boldsymbol y}). \tag{9}\]

由式（8）记

\[B_x=\frac{\mu_0m_d}{4\pi r^3}(3\cos^2\theta-1),\]

\[B_y=\frac{3\mu_0m_d}{4\pi r^3}\sin\theta\cos\theta. \tag{10}\]

因为
\(\hat{\boldsymbol z}\times\hat{\boldsymbol x}=\hat{\boldsymbol y}\)、\(\hat{\boldsymbol z}\times\hat{\boldsymbol y}=-\hat{\boldsymbol x}\)，且速度沿
\(-\hat{\boldsymbol z}\)，负电荷所受洛伦兹力分量为

先计算速度与磁场的叉积：

\[\begin{aligned}
\boldsymbol v\times\boldsymbol B
&=(-\omega'\rho\hat{\boldsymbol z})
\times(B_x\hat{\boldsymbol x}+B_y\hat{\boldsymbol y})\\
&=\omega'\rho B_y\hat{\boldsymbol x}
-\omega'\rho B_x\hat{\boldsymbol y}.
\end{aligned}\]

粒子电量为 \(-q\)，所以磁力为

\[\boldsymbol F_B=-q(\boldsymbol v\times\boldsymbol B)
=-q\omega'\rho B_y\hat{\boldsymbol x}
+q\omega'\rho B_x\hat{\boldsymbol y}.\]

再与电力 \(-q\boldsymbol E\) 相加，得到

\[F_x=-qE_x-q\omega'\rho B_y,\]

\[F_y=-qE_y+q\omega'\rho B_x.\]

轴向无加速度，故 \(F_x=0\)：

\[-\frac{qQ_0}{4\pi\varepsilon_0r^2}\cos\theta
-q\omega'r\sin\theta B_y=0. \tag{11}\]

垂直于轴的方向提供向心力：

\[-\frac{qQ_0}{4\pi\varepsilon_0r^2}\sin\theta
+q\omega'r\sin\theta B_x
=-m\omega'^2r\sin\theta. \tag{12}\]

把 \(m_d=-Q_0\omega R^2/3\)
代入式（10），再联立式（11）、（12）。先用式（11）消去库仑项，可将式（12）化为

具体地，把式（11）除以 \(-q\cos\theta\)，并代入 \(B_y\)：

\[\frac{Q_0}{4\pi\varepsilon_0r^2}
+\omega'r\sin\theta
\frac{3\mu_0m_d}{4\pi r^3}\sin\theta=0.\]

整理为

\[\frac{Q_0}{4\pi\varepsilon_0r^2}
=-\frac{3\mu_0m_d\omega'}{4\pi r^2}\sin^2\theta. \tag{12a}\]

式（12）除以 \(r\sin\theta\) 后为

\[-\frac{qQ_0}{4\pi\varepsilon_0r^3}
+q\omega'B_x=-m\omega'^2.\]

再用式（12a）替换库仑项，并代入 \(B_x\)：

\[\frac{3q\mu_0m_d\omega'}{4\pi r^3}\sin^2\theta
+\frac{q\mu_0m_d\omega'}{4\pi r^3}(3\cos^2\theta-1)
=-m\omega'^2.\]

利用 \(\sin^2\theta+\cos^2\theta=1\)，左侧括号化为

\[3\sin^2\theta+3\cos^2\theta-1=2,\]

于是

\[m\omega'= -2q\frac{\mu_0m_d}{4\pi r^3}.\]

因 \(m_d<0\)，得到正的角速度大小

\[\boxed{
\omega'=\frac{\mu_0qQ_0R^2}{6\pi mr^3}\,\omega
}. \tag{13}\]

另一方面，式（11）给出

\[\omega'\sin^2\theta=\frac{c^2}{\omega R^2}, \tag{14}\]

其中使用了 \(1/(\mu_0\varepsilon_0)=c^2\)。将式（13）代入式（14）：

\[\sin^2\theta
=\frac{c^2}{\omega^2R^2}
\frac{6\pi mr^3}{\mu_0qQ_0R^2}.\]

故

\[\boxed{
\theta=\arcsin\left[
\frac{c}{\omega R}
\sqrt{\frac{6\pi mr^3}{\mu_0qQ_0R^2}}
\right]
}. \tag{15}\]

以上先假定 \(\cos\theta\ne0\)。记式（15）根号结果为 \(u\)；若
\(0<u<1\)，两条非赤道形式解为 \(\theta=\arcsin u\) 和
\(\theta=\pi-\arcsin u\)，分别位于两侧半球，不能只保留反正弦主值。

\textbf{赤道分支。} 当 \(\theta=\pi/2\)
时，\(E_x=B_y=0\)，式（11）恒成立，不能除以 \(\cos\theta\)。记

\[b=\frac{\mu_0Q_0\omega R^2}{12\pi r^3}>0,\qquad K=\frac{Q_0}{4\pi\varepsilon_0r^3}>0.\]

在该点 \(B_x=b\)。用带符号的 \(\omega'\) 定义
\(\boldsymbol v=-\omega'r\hat{\boldsymbol z}\)，由式（12）得到

\[m\omega'^2+qb\omega'-qK=0,\]

所以赤道上有两个相反环绕方向的解

\[\boxed{\omega'_\pm=\frac{-qb\pm\sqrt{q^2b^2+4mqK}}{2m},\qquad\theta=\frac\pi2}.\]

\textbf{非赤道形式解的适用性。} 式（14）还给出轨道速率

\[v=\omega'r\sin\theta=\frac{c^2r}{\omega R^2\sin\theta}.\]

若球壳表面速率 \(\omega R<c\) 且 \(r\gg R\)，该式要求
\(v>c\)。因此非赤道解只能作为题设牛顿动力学的形式解，不能解释成满足相对论约束的实际轨道。赤道分支不受式（14）约束；使用牛顿公式仍须检查其
\(|\omega'|r\ll c\)。图 4-2
的力分解用于建立分量方程，示意箭头不代表已经满足所有轨道存在条件。

求球壳磁矩时可把球面分成纬度为 \(\theta\)、宽度为 \(R\,\mathrm d\theta\)
的窄环。该环电量、等效电流和环面积依次为

\[\begin{aligned}
\mathrm dq&=\sigma(2\pi R\sin\theta)(R\,\mathrm d\theta),\\
\mathrm dI&=\frac{\omega\,\mathrm dq}{2\pi},\\
S_\theta&=\pi R^2\sin^2\theta.
\end{aligned}\]

三者相乘产生 \(\sin^3\theta\)，积分方向由 \(\boldsymbol\omega\) 的负
\(x\) 方向确定。这样推导比直接套均匀带电球壳磁矩公式更能控制符号。

远区场只保留磁偶极项，其 \(B_r\)、\(B_\theta\)
形式可由电偶极场逐项对应，但要注意磁矩本身带方向。研究负电荷圆周运动时，把洛伦兹力分解成沿
\(x\)
轴和指向圆周轴心的两个分量：轴向加速度为零给出一条方程，横向分量提供
\(m\omega'^2r\sin\theta\)。先用轴向方程消去一个量，再求角速度，能避免在符号上混乱。

结果需要两层存在性检查：根号内为正决定角速度实数，\(\sin\theta\leqslant1\)
决定所给参数下是否存在这种定常纬圈。改变电荷符号或转动方向时，洛伦兹力方向改变，允许的环绕方向也随之改变，不能只取角速度的绝对值而忽略方向。

\newpage

\section{第 5 题}\label{ux7b2c-5-ux9898}

如图（a），光滑水平面上有两个均匀绝缘圆环，半径分别为 \(2R\) 和
\(R\)，质量均为 \(m\)。轻杆 \(O_1O_2\)
在两圆环的圆心处分别与圆环铰接，两圆环接触处始终无相对滑动。在小圆环
\(O_2\) 的右端固定电量为 \(q\) 的轻质点电荷。

\begin{enumerate}
\def\labelenumi{\arabic{enumi}.}
\tightlist
\item
  在以 \(O_1\) 为圆心、半径为 \(2R\) 的区域内加竖直方向的均匀磁场，且
  \(B=B_0+kt\)。固定圆环 \(O_1\)，将系统从图（a）位置静止释放。求杆
  \(O_1O_2\) 转过角度 \(\theta\) 时的角速度 \(\omega_1\) 和角加速度
  \(\beta_1\)。
\item
  接第（1）问，但允许圆环 \(O_1\) 绕过自身圆心的竖直轴自由转动。已知杆
  \(O_1O_2\) 转过角度 \(\theta\) 的过程用时为 \(t\)，求杆的角速度
  \(\omega_2\) 和角加速度 \(\beta_2\)。允许用点电荷相对 \(O_1\) 的方位角
  \(\theta'\) 及 \(\dot\theta'\) 表示结果，但必须给出 \(\theta'\)
  的表达式。
\item
  如图（b），固定圆环 \(O_1\)，将 \(O_2\) 换成半径为 \(2R\)、质量仍为
  \(m\) 的圆环 \(O_3\)。磁场区域扩展为以 \(O_1\) 为圆心、半径大于 \(6R\)
  的圆形区域，仍有 \(B=B_0+kt\)，两环接触处无相对滑动，点电荷仍为
  \(q\)。求系统从图（b）位置静止释放瞬间杆 \(O_1O_3\) 的角加速度
  \(\beta_3\)。
\item
  接第（3）问，求杆 \(O_1O_3\) 转过 \(180^\circ\) 时的角速度
  \(\omega_4\) 和角加速度 \(\beta_4\)。
\end{enumerate}

\begin{figure}[H]
\centering
\includegraphics[width=0.9\linewidth,height=\textheight,keepaspectratio,alt={图 5-1　两种外切圆环系统：（a）半径分别为 2R、R；（b）两环半径均为 2R。}]{../figures/question-05.pdf}
\caption*{图 5-1　两种外切圆环系统：（a）半径分别为
2R、R；（b）两环半径均为 2R。}
\end{figure}

\subsection{一、题型定位}\label{ux4e00ux9898ux578bux5b9aux4f4d-4}

由两圆环的接触约束求角速度，再用法拉第定律确定感生电场。动力学计算要同时计入圆心公转与圆环自转，第二问还需处理另一圆环的转动。

\subsection{二、知识点}\label{ux4e8cux77e5ux8bc6ux70b9-4}

\begin{itemize}
\tightlist
\item
  轴对称变化磁场产生的环形感生电场。
\item
  外切圆无滑动的接触点速度约束。
\item
  点电荷位置矢量的复数或直角坐标表示。
\item
  功---能定理与角动量定理。
\item
  薄圆环转动惯量 \(I=mR^2\)。
\item
  \(\operatorname{atan2}\) 方位角及其导数。
\end{itemize}

\subsection{三、物理图像}\label{ux4e09ux7269ux7406ux56feux50cf-4}

固定大环时，小环圆心绕 \(O_1\)
公转，同时小环为满足接触点无滑动而自转。点电荷并不沿以 \(O_1\)
为圆心的圆运动，它的位置是``圆心矢量 +
环上半径矢量''的和。因此感生电场做功应写成
\(qE_\varphi r\,\mathrm d\theta'\)，不能直接用杆角 \(\theta\) 代替
\(\theta'\)。

\subsection{四、解题入口}\label{ux56dbux89e3ux9898ux5165ux53e3-4}

所有小问先做两步：第一步由接触点速度相等求圆环自转角；第二步写点电荷相对
\(O_1\) 的位置矢量，求其方位角 \(\theta'\)。随后再根据磁场区域判断

\[E_\varphi=\frac{2kR^2}{r}\quad\text{或}\quad E_\varphi=\frac{kr}{2}.\]

前者适用于点电荷位于半径 \(2R\)
的变化磁通区域外，后者适用于点电荷处在均匀变化磁场内部。

\subsection{五、解析}\label{ux4e94ux89e3ux6790-4}

图 5-2 左侧标出两外切圆环的公转角 \(\theta\)、小环自转角
\(\varphi\)、点电荷位置 \(\bm r_q=3R\bm e(\theta)+R\bm e(\varphi)\)
以及感生电场方向；右侧区分点电荷位于变化磁通区域外、内时的环向电场公式。第（1）---（4）步均先由无滑动确定角度关系，再按电荷实际位置选取电场表达式计算功和动力学。

\begin{figure}[H]
\centering
\includegraphics[width=0.9\linewidth,height=\textheight,keepaspectratio,alt={图 5-2　外切圆环的滚动角关系、点电荷位置矢量和感生电场做功路径。}]{../figures/solution-05.pdf}
\caption*{图
5-2　外切圆环的滚动角关系、点电荷位置矢量和感生电场做功路径。}
\end{figure}

约定 \(qk>0\) 时，系统沿图中 \(\theta\) 增大的方向运动；若
\(qk<0\)，所有带方向的角速度、角加速度同时反向。

\textbf{（1）固定 \(O_1\)，小环半径为 \(R\)}

按图 5-2 左侧，小环圆心以半径 \(3R\)
公转，而接触点相对小环圆心的自转半径为 \(R\)。接触点速度为零给出
\(3R\dot\theta-R\dot\varphi=0\)，这是后续位置矢量和感生电场做功的几何入口。

小环圆心 \(O_2\) 到 \(O_1\) 的距离为 \(3R\)。设小环相对实验室的自转角为
\(\varphi\)。接触点无相对滑动，而大环固定，所以小环接触点的速度为零：

\[3R\dot\theta-R\dot\varphi=0.\]

由初始时 \(\theta=\varphi=0\)，得到

\[\boxed{\varphi=3\theta}. \tag{1}\]

点电荷相对 \(O_1\) 的位置为

\[\boldsymbol r_q
=3R(\cos\theta,\sin\theta)
+R(\cos3\theta,\sin3\theta). \tag{2}\]

故其方位角 \(\theta'\) 应用连续的 \(\operatorname{atan2}\) 表示：

\[\boxed{
\theta'=\operatorname{atan2}
\left(3\sin\theta+\sin3\theta,
3\cos\theta+\cos3\theta\right)
}. \tag{3}\]

变化磁场只存在于半径 \(2R\) 的区域内。点电荷在该区域外，法拉第定律给出

\[E_\varphi(2\pi r_q)=\pi(2R)^2k,\]

即

\[E_\varphi=\frac{2kR^2}{r_q}. \tag{4}\]

点电荷转过微小方位角 \(\mathrm d\theta'\) 时，感生电场做功

\[\mathrm dW=qE_\varphi r_q\,\mathrm d\theta'
=2kqR^2\,\mathrm d\theta'.\]

从静止释放到当前位置，

\[W=2kqR^2\theta'. \tag{5}\]

小环的平动动能与自转动能分别为

\[K_{\rm tr}=\frac12m(3R\dot\theta)^2,\]

\[K_{\rm rot}=\frac12(mR^2)(3\dot\theta)^2.\]

所以

\[K=9mR^2\dot\theta^2. \tag{6}\]

由功---能定理 \(K=W\)：

\[9mR^2\dot\theta^2=2kqR^2\theta'.\]

因此

\[\boxed{
\omega_1=\dot\theta
=\sqrt{\frac{2kq}{9m}\,\theta'}
}. \tag{7}\]

为求角加速度，对式（6）与式（5）关于时间求导：

\[18mR^2\dot\theta\ddot\theta
=2kqR^2\dot\theta'.\]

又 \(\dot\theta'=(\mathrm d\theta'/\mathrm d\theta)\dot\theta\)，约去
\(\dot\theta\)：

\[\ddot\theta=\frac{kq}{9m}\frac{\mathrm d\theta'}{\mathrm d\theta}. \tag{8}\]

由式（3）求导可得

\[\frac{\mathrm d\theta'}{\mathrm d\theta}
=\frac{6(1+\cos2\theta)}{5+3\cos2\theta}. \tag{9}\]

所以

\[\boxed{
\beta_1=\ddot\theta
=\frac{2kq(1+\cos2\theta)}{3m(5+3\cos2\theta)}
}. \tag{10}\]

当 \(\theta=0\) 时，式（10）给出
\(\beta_1=kq/(3m)\)，可作为释放瞬间检查。

\textbf{（2）大环 \(O_1\) 可自由转动}

设大环自转角为 \(\alpha\)，小环自转角仍为 \(\varphi\)。接触点速度相等：

\[2R\dot\alpha=3R\dot\theta-R\dot\varphi.\]

积分并用初始角均为零：

\[\boxed{\varphi+2\alpha=3\theta}. \tag{11}\]

感生电场对整个机械系统关于 \(O_1\) 的外力矩为

\[\tau=qE_\varphi r_q=2kqR^2. \tag{12}\]

系统关于 \(O_1\) 的角动量包括大环自转、小环圆心公转和小环自转：

\[L=4mR^2\dot\alpha+9mR^2\dot\theta+mR^2\dot\varphi. \tag{13}\]

由 \(L=\tau t\) 及 \(\dot\varphi=3\dot\theta-2\dot\alpha\)：

\[2\dot\alpha+12\dot\theta=\frac{2kq}{m}t.\]

定义

\[a=\frac{kq}{m},\qquad \Omega_0=at,\]

则

\[\boxed{\dot\alpha=-6\dot\theta+\Omega_0},\]

\[\boxed{\dot\varphi=15\dot\theta-2\Omega_0}. \tag{14}\]

进一步积分：

\[\alpha=-6\theta+\frac12at^2,\]

\[\varphi=15\theta-at^2=15\theta-\Omega_0t. \tag{15}\]

因此点电荷方位角为

\[\boxed{
\theta'=\operatorname{atan2}
\left[3\sin\theta+\sin(15\theta-at^2),
3\cos\theta+\cos(15\theta-at^2)
\right]
}. \tag{16}\]

总动能为

\[K=\frac12mR^2
\left(4\dot\alpha^2+9\dot\theta^2+\dot\varphi^2\right). \tag{17}\]

将式（14）代入：

\[K=\frac12mR^2
\left(378\dot\theta^2-108\Omega_0\dot\theta+8\Omega_0^2\right). \tag{18}\]

仍有 \(W=2kqR^2\theta'\)。令 \(K=W\)，并除以 \(mR^2/2\)：

\[378\dot\theta^2-108\Omega_0\dot\theta+8\Omega_0^2
=4a\theta'. \tag{19}\]

解关于 \(\dot\theta\) 的二次方程：

\[\boxed{
\omega_2=\dot\theta
=\frac1{63}\left[
9\Omega_0\pm\sqrt{3(14a\theta'-\Omega_0^2)}
\right]
}. \tag{20}\]

两支都在 \(t=0\)
汇合为零，必须用释放瞬间的动力学选支。在下文的正方向约定 \(qk>0\)
下，拉格朗日方程给出
\(\ddot\theta(0)=a/6\)，对应式（20）的正号支；若采用相反的感生电场方向，所有角量反向并选择相应连续分支。

对式（20）求导。记

\[D=3(14a\theta'-\Omega_0^2),\qquad \dot\Omega_0=a,\]

则

\[\boxed{
\beta_2=\ddot\theta
=\frac1{63}\left[
9a\pm\frac{3a(7\dot\theta'-\Omega_0)}{\sqrt D}
\right]
}. \tag{21}\]

式（21）中的正负号必须与式（20）取同一分支；\(t\to0\)
时应以极限求值，不能直接代入造成 \(0/0\)。

\textbf{（3）固定 \(O_1\)，两环半径均为 \(2R\)}

本问磁场区域的半径大于 \(6R\)，而点电荷到 \(O_1\) 的距离不超过
\(6R\)，所以全程使用区内公式 \(E_\varphi=kr_q/2\)。图 5-2
右图先给出一般分区关系；这里无需分段。

两圆心距离为 \(4R\)。无滑动条件为

\[4R\dot\theta-2R\dot\varphi=0,\]

故

\[\boxed{\varphi=2\theta}. \tag{22}\]

点电荷位置

\[\boldsymbol r_q=4R\boldsymbol e(\theta)+2R\boldsymbol e(2\theta),\]

所以

\[r_q^2=R^2(20+16\cos\theta), \tag{23}\]

\[\theta'=\operatorname{atan2}
\left(2\sin\theta+\sin2\theta,
2\cos\theta+\cos2\theta\right). \tag{24}\]

此时点电荷始终位于均匀变化磁场内，故

\[E_\varphi=\frac{kr_q}{2}.\]

微元功为

\[\begin{aligned}
\mathrm dW
&=q\frac{kr_q}{2}r_q\,\mathrm d\theta'\\
&=\frac{kq}{2}r_q^2\frac{\mathrm d\theta'}{\mathrm d\theta}\,\mathrm d\theta\\
&=16kqR^2(1+\cos\theta)\,\mathrm d\theta.
\end{aligned} \tag{25}\]

系统动能为

\[\begin{aligned}
K
&=\frac12m(4R\dot\theta)^2
+\frac12m(2R)^2(2\dot\theta)^2\\
&=16mR^2\dot\theta^2.
\end{aligned} \tag{26}\]

由 \(\mathrm dK=\mathrm dW\)：

\[32mR^2\dot\theta\ddot\theta\,\mathrm dt
=16kqR^2(1+\cos\theta)\dot\theta\,\mathrm dt.\]

于是

\[\boxed{
\ddot\theta=\frac{kq}{2m}(1+\cos\theta)
}. \tag{27}\]

释放瞬间 \(\theta=0\)，所以

\[\boxed{\beta_3=\frac{kq}{m}}. \tag{28}\]

\textbf{（4）杆转过 \(180^\circ\)}

用 \(\ddot\theta=\dot\theta\,\mathrm d\dot\theta/\mathrm d\theta\)
积分：

\[\dot\theta\,\mathrm d\dot\theta
=\frac{kq}{2m}(1+\cos\theta)\,\mathrm d\theta.\]

从 \(0\) 积到 \(\theta\)：

\[\frac12\dot\theta^2
=\frac{kq}{2m}(\theta+\sin\theta),\]

因此

\[\boxed{
\dot\theta^2=\frac{kq}{m}(\theta+\sin\theta)
}. \tag{29}\]

当 \(\theta=\pi\) 时，\(\sin\pi=0\)，故

\[\boxed{\omega_4=\sqrt{\frac{\pi kq}{m}}}. \tag{30}\]

同时由式（27），

\[\boxed{\beta_4=0}. \tag{31}\]

代入 \(\theta=\pi\) 时必须保留弧度值 \(\pi\)，不能把它误写成数值
\(2\)；式（30）由式（29）直接代入得到。

本题应先把滚动几何完全解决，再进入动力学。外切圆环无相对滑动意味着接触点速度相等，杆的公转角、每个圆环的自转角以及点电荷绕
\(O_1\)
的方位角并不相同。第（1）问固定大环时，点电荷的方位角可由``小环圆心公转
+ 小环相对杆自转''合成；感生电场沿以 \(O_1\) 为圆心的圆周方向，功率为
\(qE_\varphi\rho\dot\theta'\)。

由于 \(B=B_0+kt\)，真正产生感生电场的是 \(\dot B=k\)，常量 \(B_0\)
不进入法拉第电场。按题图取磁场入纸面为正、逆时针方位角为正，法拉第定律给出
\(2\pi\rho E_\varphi=\pi\rho^2k\)，即
\(E_\varphi=k\rho/2\)；若磁场取出纸面为正，该分量整体反号。再通过功---能定理把电场功转化为各部分平动、转动动能。计算等效转动惯量时，每个圆环既有圆心平动，也有绕心自转，不能只取杆端的质点模型。

第（2）问大环可自由转动，外界关于竖直轴的力矩来自感生电场对点电荷的作用。角动量定理和功---能关系分别提供关于广义速度的线性、二次约束；二次方程出现两个代数根时，应从
\(t=0\)
的角加速度方向选取与初态连续的分支。第（3）、（4）问两环等半径，几何关系改变，最终在
\(\theta=\pi\) 代入时应保留 \(\pi\) 而不是用行程长度的数值替代角度。

\newpage

\section{第 6 题}\label{ux7b2c-6-ux9898}

一容器中装有质量 \(m_w=9\,\mathrm g\) 的水和物质的量
\(n_a=0.5\,\mathrm{mol}\) 的空气。初始时容器体积为
\(V_0=61.4\,\mathrm L\)。随后绝热压缩到体积
\(V_1\)，并认为该过程中水全部处于气态。压缩后气体温度为
\(T_1=373\,\mathrm K\)、总压强为 \(p_1=1\,\mathrm{atm}\)，其中
\(1\,\mathrm{atm}=1.01\times10^5\,\mathrm{Pa}\)。再使系统在 \(T_1\)
下等温压缩到体积 \(V_1/4\)，并认为系统中有足够的凝结核。

本题条件下，水的汽化热为
\(L=2.26\times10^5\,\mathrm{J\,kg^{-1}}\)，液态水体积可忽略。气体常量为

\[R=8.31\,\mathrm{J\,mol^{-1}\,K^{-1}}.\]

水蒸气和空气分别视为多原子分子、双原子分子理想气体，只有平动和转动自由度被激发。

\begin{enumerate}
\def\labelenumi{\arabic{enumi}.}
\tightlist
\item
  求绝热压缩前水蒸气与空气的分压 \(p_{w0}\)、\(p_{a0}\)。
\item
  求随后等温压缩过程中系统放出的热量。
\end{enumerate}

\subsection{一、题型定位}\label{ux4e00ux9898ux578bux5b9aux4f4d-5}

空气与水蒸气混合物先绝热压缩，后等温压缩。前一段用混合热容求绝热指数，后一段按未饱和、饱和起点和凝结阶段分段积分。

\subsection{二、知识点}\label{ux4e8cux77e5ux8bc6ux70b9-5}

\begin{itemize}
\tightlist
\item
  道尔顿分压定律。
\item
  分子自由度与定容摩尔热容。
\item
  混合理想气体的总热容和绝热指数。
\item
  理想气体绝热关系 \(pV^\gamma=\text{常量}\)。
\item
  等温压缩功与相变潜热。
\item
  饱和蒸气压约束和凝结质量。
\end{itemize}

\subsection{三、物理图像}\label{ux4e09ux7269ux7406ux56feux50cf-5}

在 \(373\,\mathrm K\)、\(p_1=1\,\mathrm{atm}\)
的状态，空气和水蒸气物质的量相等，所以两者分压均为
\(0.5\,\mathrm{atm}\)。水在 \(373\,\mathrm K\) 的饱和蒸气压为
\(1\,\mathrm{atm}\)，开始时仍未饱和；等温压缩到 \(V_1/2\)
时才刚好饱和，此后继续压缩会有水凝结，水蒸气分压保持
\(1\,\mathrm{atm}\)。

\subsection{四、解题入口}\label{ux56dbux89e3ux9898ux5165ux53e3-5}

先求水的物质的量，确认它与空气相等。再由压缩后状态求 \(V_1\)，得到
\(V_1/V_0=1/2\)。绝热阶段对混合物使用总
\(C_V\)；等温阶段把空气和水蒸气分开记账，水蒸气还要再加凝结潜热。

\subsection{五、解析}\label{ux4e94ux89e3ux6790-5}

图 6-1 画出 \(373\,\mathrm K\) 等温压缩时水蒸气分压 \(p_w(V)\)：从
\(V_1\) 压到 \(V_1/2\) 前尚未饱和，满足 \(p_wV=n_wRT\)；达到饱和蒸气压
\(p_s=p_1\) 后继续压到
\(V_1/4\)，分压保持不变并有水蒸气凝结。第（2）问必须按这个转折点分两段计算热量。

\begin{figure}[H]
\centering
\includegraphics[width=0.9\linewidth,height=\textheight,keepaspectratio,alt={图 6-1　等温压缩中水蒸气分压随体积变化：先按 1/V 增大，达到饱和后保持不变并发生凝结。}]{../figures/solution-06.pdf}
\caption*{图 6-1　等温压缩中水蒸气分压随体积变化：先按 \(1/V\)
增大，达到饱和后保持不变并发生凝结。}
\end{figure}

水的摩尔质量为 \(M_w=18\,\mathrm{g\,mol^{-1}}\)，所以

\[n_w=\frac{m_w}{M_w}
=\frac{9\,\mathrm g}{18\,\mathrm{g\,mol^{-1}}}
=0.5\,\mathrm{mol}. \tag{1}\]

可见

\[n_w=n_a=0.5\,\mathrm{mol},\qquad n=n_w+n_a=1.0\,\mathrm{mol}. \tag{2}\]

压缩后由理想气体状态方程

\[p_1V_1=nRT_1,\]

得到

\[\begin{aligned}
V_1
&=\frac{1.0\times8.31\times373}{1.01\times10^5}\,\mathrm{m^3}\\
&=3.07\times10^{-2}\,\mathrm{m^3}
\approx30.7\,\mathrm L.
\end{aligned} \tag{3}\]

因此题给 \(V_0=61.4\,\mathrm L\) 恰为

\[\frac{V_1}{V_0}=\frac12. \tag{4}\]

这个体积比应先单独保留，因为后面所有压强比和饱和节点都由它导出。数值上还可反算检查：

\[p_1V_1
=(1.01\times10^5)(3.07\times10^{-2})\,\mathrm J
\approx3.10\times10^3\,\mathrm J,\]

而

\[nRT_1
=(1.0)(8.31)(373)\,\mathrm J
\approx3.10\times10^3\,\mathrm J.\]

两者一致，说明升、毫升与立方米之间的换算没有出错。

\textbf{（1）绝热压缩前的分压}

水蒸气视为非线性多原子分子，平动和转动自由度共 \(6\) 个：

\[C_{V,w}=3R.\]

空气视为双原子分子，平动和转动自由度共 \(5\) 个：

\[C_{V,a}=\frac52R.\]

混合物的总定容热容为

\[C_V^{\rm tot}=n_wC_{V,w}+n_aC_{V,a}
=\frac12(3R)+\frac12\left(\frac52R\right)
=\frac{11}{4}R. \tag{5}\]

总物质的量为 \(1\,\mathrm{mol}\)，故总定压热容

\[C_P^{\rm tot}=C_V^{\rm tot}+R=\frac{15}{4}R.\]

等效绝热指数

\[\boxed{\gamma=\frac{C_P^{\rm tot}}{C_V^{\rm tot}}=\frac{15}{11}}. \tag{6}\]

这里不能把两个组分的绝热指数直接作算术平均。例如水蒸气与空气各占
\(0.5\,\mathrm{mol}\)，错误做法会得到

\[\frac12\left(\frac43+\frac75\right),\]

但绝热过程真正相加的是内能和焓，因此必须先形成

\[C_V^{\rm tot}=\sum_i n_iC_{V,i},\]

\[C_P^{\rm tot}=C_V^{\rm tot}+nR,\]

最后才取二者之比。式（6）正是这一热容加权过程的结果。

绝热关系为

\[p_0V_0^\gamma=p_1V_1^\gamma.\]

于是

\[p_0=p_1\left(\frac{V_1}{V_0}\right)^\gamma
=p_1\,2^{-15/11}. \tag{7}\]

初态两种气体同处 \(T_0,V_0\)，分压与物质的量成正比；因两者物质的量相等，

\[p_{w0}=p_{a0}=\frac12p_0.\]

所以

\[\boxed{
p_{w0}=p_{a0}
=\frac{p_1}{2}\,2^{-15/11}
\approx1.96\times10^4\,\mathrm{Pa}
}. \tag{8}\]

\textbf{（2）等温压缩放热量}

沿图 6-1 的橙色箭头，第一段从 \(V_1\) 到 \(V_1/2\)
是空气和未饱和水蒸气共同等温压缩；第二段从 \(V_1/2\) 到 \(V_1/4\)
时水蒸气分压钉在 \(p_s\)，凝结放热用汽化潜热 \(L\) 计算。这里 \(L\)
是两相焓差，已包括相变对应的压力体积功，不再另加一份水蒸气的
\(p_s\Delta V\)。

在状态 1，两气体分压相等：

\[p_{w1}=p_{a1}=\frac12p_1=0.5\,\mathrm{atm}. \tag{9}\]

\(373\,\mathrm K\) 时水的饱和蒸气压为
\(p_s=1\,\mathrm{atm}=p_1\)。在尚未凝结时，\(p_wV=n_wRT_1\)，故水蒸气分压从
\(0.5\,\mathrm{atm}\) 增至 \(1\,\mathrm{atm}\) 时体积减半：

\[V_s=\frac12V_1. \tag{10}\]

也可以直接用水蒸气的物质的量核验：

\[V_s=\frac{n_wRT_1}{p_s}
=\frac{0.5RT_1}{p_1}.\]

又因为

\[V_1=\frac{(n_a+n_w)RT_1}{p_1}
=\frac{1.0RT_1}{p_1},\]

故 \(V_s/V_1=0.5\)。这说明``分压由 \(0.5\,\mathrm{atm}\) 加倍到
\(1\,\mathrm{atm}\)''和``体积减半''是同一件事的两种表述。

空气在整个等温过程中始终不凝结，从 \(V_1\) 压缩到
\(V_1/4\)，放热量等于对空气所做的功：

\[\begin{aligned}
Q_a
&=n_aRT_1\ln\frac{V_1}{V_1/4}\\
&=n_aRT_1\ln4.
\end{aligned} \tag{11}\]

水蒸气在 \(V_1\to V_s=V_1/2\) 的未饱和等温压缩阶段放热

\[Q_{w,\rm gas}=n_wRT_1\ln\frac{V_1}{V_s}
=n_wRT_1\ln2. \tag{12}\]

从 \(V_s\) 继续压到 \(V_f=V_1/4\) 时，水蒸气分压保持
\(p_s=p_1\)。末态水蒸气物质的量为

\[n_{w,f}=\frac{p_sV_f}{RT_1}
=\frac{p_1(V_1/4)}{RT_1}.\]

又 \(p_1V_1=(n_a+n_w)RT_1=1\,\mathrm{mol}\cdot RT_1\)，所以

\[n_{w,f}=0.25\,\mathrm{mol}. \tag{13}\]

从饱和节点到末态，水蒸气压与温度都保持不变，因此气相物质的量与体积成正比：

\[\frac{n_{w,f}}{n_{w,s}}
=\frac{V_f}{V_s}
=\frac{V_1/4}{V_1/2}
=\frac12.\]

在饱和节点尚未发生宏观凝结，\(n_{w,s}=0.5\,\mathrm{mol}\)，所以同样得到

\[n_{w,f}=0.25\,\mathrm{mol}.\]

这一比例法不需要再次代入气体常数，可作为式（13）的独立检查。

凝结的水为

\[\Delta n_w=0.50-0.25=0.25\,\mathrm{mol},\]

\[\Delta m=\Delta n_wM_w
=0.25\times18\,\mathrm g=4.5\,\mathrm g. \tag{14}\]

凝结放出的潜热为

\[Q_L=\Delta mL
=4.5\times10^{-3}\times2.26\times10^5
=1.017\times10^3\,\mathrm J. \tag{15}\]

总放热量

\[\begin{aligned}
Q
&=Q_a+Q_{w,\rm gas}+Q_L\\
&=0.5RT_1\ln4+0.5RT_1\ln2+\Delta mL\\
&=\frac32RT_1\ln2+\Delta mL.
\end{aligned}\]

代入数值：

\[\boxed{Q\approx4.24\times10^3\,\mathrm J=4.24\,\mathrm{kJ}}. \tag{16}\]

三部分热量的数量级也可以分别列出：

\[Q_a=\frac12RT_1\ln4\approx2.15\,\mathrm{kJ},\]

\[Q_{w,\rm gas}=\frac12RT_1\ln2\approx1.07\,\mathrm{kJ},\]

\[Q_L\approx1.02\,\mathrm{kJ}.\]

相加得到

\[Q\approx2.15+1.07+1.02=4.24\,\mathrm{kJ}.\]

这里报告的是``向外放出的热量大小''。若统一采用 \(\delta Q>0\)
表示系统吸热，则整个等温压缩过程应写成

\[Q_{\rm sys}=-4.24\,\mathrm{kJ}.\]

把文字约定与代数符号分开写，可以避免同一个正数既被解释为吸热又被解释为放热。

初态虽同时含空气和水蒸气，但两者占据同一体积、温度相同，因此分别使用理想气体状态方程即可。由水质量先换算物质的量
\(n_w=m_w/M_w=0.5\,\mathrm{mol}\)，恰与空气物质的量相等，所以初态两种气体分压相等。绝热压缩末态给出总压强和温度，可由总物质的量求
\(V_1\)；这一数值还要与后续饱和条件一起使用。

等温压缩应先找``开始凝结''的体积，而不能从一开始就加入汽化热。初始等温阶段两种气体的物质的量固定，分压均按
\(1/V\) 增大；当水蒸气分压达到 \(373\,\mathrm K\)
下的饱和蒸气压后，继续压缩时水蒸气分压保持不变，多余水蒸气凝结。空气则始终不凝结，仍满足
\(p_aV=n_aRT\)。

放热量分三部分相加：空气在全部体积变化中的等温压缩热、水蒸气在未凝结阶段的等温压缩热，以及凝结质量乘汽化热。凝结后的水蒸气不能再使用固定
\(n\) 的 \(nRT\ln(V_i/V_f)\) 跨越整个区间。最后检查 \(Q>0\)
表示``系统放热''的大小；若采用热力学中系统吸热为正的约定，则系统的
\(Q_{\rm sys}\) 应写成负值。

\newpage

\section{第 7 题}\label{ux7b2c-7-ux9898}

在惯性系 \(S\) 中有一半径为 \(R\) 的圆 \(C\)，圆周上有一点 \(O\)。在
\(t=0\) 时，质点 \(A\)、\(B\) 同时从 \(O\) 出发，均以恒定速率
\(\beta c\) 运动。\(A\) 沿圆周运动，\(B\) 沿圆在 \(O\)
点的切线做匀速直线运动，初始时二者速度同向。

\(A\)、\(B\) 各携带一只钟。\(B\) 的钟指示 \(B\) 参考系中的时间，并在
\(O\) 点调为零；\(A\) 的钟与静止在 \(S\) 系、指示 \(S\) 系时间的钟在
\(S\) 系中同步。

\begin{enumerate}
\def\labelenumi{\arabic{enumi}.}
\tightlist
\item
  求当 \(B\) 认为自己的钟指示 \(t_B\) 时，\(B\) 认为 \(A\) 的钟示数
  \(t_A\) 所满足的方程，不必求解。
\item
  求当 \(A\) 认为自己的钟指示 \(t_A\) 时，\(A\) 在其瞬时惯性系中认为
  \(B\) 的钟示数 \(t_B\)。
\end{enumerate}

\begin{figure}[H]
\centering
\includegraphics[width=0.844\linewidth,height=\textheight,keepaspectratio,alt={图 7-1　A 沿圆周运动，B 沿 O 点切线运动，二者初速度同向。}]{../figures/question-07.pdf}
\caption*{图 7-1　A 沿圆周运动，B 沿 O 点切线运动，二者初速度同向。}
\end{figure}

\subsection{一、题型定位}\label{ux4e00ux9898ux578bux5b9aux4f4d-6}

在指定事件处比较两个运动者的同时性。惯性运动者可用一套洛伦兹变换；圆周运动者需在该事件使用瞬时共动系，不能把它推广到整条轨迹。

\subsection{二、知识点}\label{ux4e8cux77e5ux8bc6ux70b9-6}

\begin{itemize}
\tightlist
\item
  洛伦兹变换中的时间坐标。
\item
  运动钟示数与惯性系坐标时间的关系。
\item
  圆周运动的位置、速度参数方程。
\item
  瞬时惯性系的同时性条件。
\item
  四维间隔中``时间差减去速度方向投影''的结构。
\end{itemize}

\subsection{三、物理图像}\label{ux4e09ux7269ux7406ux56feux50cf-6}

在 \(S\) 系中，\(B\) 沿 \(x\) 轴直线运动，\(A\) 沿圆周运动。第一问是在
\(B\) 的惯性系中取与 \(B\) 事件同时的 \(A\) 事件；第二问是在 \(A\)
的瞬时共动系中，过 \(A\) 事件作一条``同时面''，求它与 \(B\)
世界线的交点。

\subsection{四、解题入口}\label{ux56dbux89e3ux9898ux5165ux53e3-6}

把 \(O\) 取为坐标原点，\(B\) 的初速度沿 \(+x\)，圆心位于 \((0,R)\)。定义

\[\Omega=\frac{\beta c}{R},\qquad \gamma=\frac1{\sqrt{1-\beta^2}}.\]

则 \(A\) 在 \(S\) 系中的位置和速度方向都可由角度 \(\Omega t\)
表示。第二问不必写完整的平移---旋转---洛伦兹矩阵，只需使用瞬时同时条件

\[c\Delta t-\boldsymbol\beta_A\cdot\Delta\boldsymbol r=0.\]

\subsection{五、解析}\label{ux4e94ux89e3ux6790-6}

图 7-2 左侧画出实验室系 \(S\) 中圆周运动者 \(A\) 与沿切线运动者 \(B\)
的空间轨迹；右侧在时空图中画出两条世界线，并用过给定 \(A\)
事件的瞬时同时面与 \(B\) 世界线相交。第（1）问从给定 \(B\) 事件用 \(B\)
系同时条件寻找 \(A\) 事件，第（2）问反过来从给定 \(A\)
事件使用其瞬时惯性系同时条件寻找 \(B\) 事件。

\begin{figure}[H]
\centering
\includegraphics[width=0.641\linewidth,height=\textheight,keepaspectratio,alt={图 7-2　S 系中的两条轨迹，以及 A 瞬时惯性系同时面与 B 世界线的交点。}]{../figures/solution-07.pdf}
\caption*{图 7-2　\(S\) 系中的两条轨迹，以及 \(A\) 瞬时惯性系同时面与
\(B\) 世界线的交点。}
\end{figure}

在 \(S\) 系中，\(A\) 的参数方程为

\[\boxed{
x_A(t)=R\sin\Omega t,\qquad
y_A(t)=R(1-\cos\Omega t)
}. \tag{1}\]

\(B\) 的世界线为

\[x_B(t)=\beta ct,\qquad y_B(t)=0. \tag{2}\]

题目规定 \(A\) 的钟与 \(S\) 系坐标时钟同步，因此这里的 \(t_A\) 就是对应
\(A\) 事件的 \(S\) 系时间；\(B\) 的钟是其固有时，所以 \(S\) 系中对应的
\(B\) 事件时间为 \(\gamma t_B\)。

两种时间的对应关系可先写成事件表。\(A\) 钟读数为 \(t_A\) 的事件在 \(S\)
系中是

\[P_A=\bigl(ct_A,\,R\sin\Omega t_A,\,R(1-\cos\Omega t_A)\bigr).\]

\(B\) 钟读数为 \(t_B\) 的事件在 \(S\) 系中是

\[P_B=\bigl(c\gamma t_B,\,\beta c\gamma t_B,\,0\bigr).\]

第一问从给定的 \(P_B\) 出发，寻找与它在 \(B\) 系同时的
\(P_A\)；第二问从给定的 \(P_A\) 出发，寻找落在 \(A\)
瞬时同时面上的另一件 \(B\)
事件。先分清``给定哪一个事件、寻找哪一个事件''，可以防止把两问错误地看成互求反函数。

\textbf{（1）\(B\) 系认为同时的 \(A\) 钟示数}

图 7-2 的时空图提醒我们，``同时''不是在 \(S\)
系中取同一水平线。对第一问应使用 \(B\) 的固定惯性系洛伦兹时间
\(t'\)，令目标 \(A\) 事件与给定 \(B\) 事件具有相同的 \(t'\)。

从 \(S\) 系到 \(B\) 系的洛伦兹时间变换为

\[ct'=\gamma(ct-\beta x). \tag{3}\]

对 \(A\) 事件，\(t=t_A\)、\(x=x_A(t_A)=R\sin(\Omega t_A)\)。\(B\)
认为该事件与自己钟读数 \(t_B\) 的事件同时，即 \(t'=t_B\)。代入式（3）：

\[ct_B=\gamma\left[ct_A-\beta R\sin(\Omega t_A)\right].\]

另一方面，把 \(B\) 自己的事件代入同一变换也可核对时标：

\[\begin{aligned}
ct'_B
&=\gamma\left(c\gamma t_B-\beta\cdot\beta c\gamma t_B\right)\\
&=\gamma^2ct_B(1-\beta^2)\\
&=ct_B.
\end{aligned}\]

所以式（3）右端确实应与 \(ct_B\) 比较，而不是与 \(c\gamma t_B\) 比较。

故所求隐式方程为

\[\boxed{
ct_B=\gamma ct_A-\gamma\beta R
\sin\left(\frac{\beta ct_A}{R}\right)
}. \tag{4}\]

该方程一般不能用初等函数显式解出 \(t_A(t_B)\)，题目也只要求列方程。

\textbf{（2）\(A\) 的瞬时惯性系认为同时的 \(B\) 钟示数}

第二问中 \(A\)
的速度方向随圆周运动不断改变，所以同时面也随事件改变。按图 7-2，过给定
\(A\) 事件作与其瞬时四速度正交的同时面，再求它与 \(B\)
世界线的交点；这不是第一问方程的简单反函数。

考虑 \(A\) 钟示数为 \(t_A\) 的事件。记

\[\psi=\Omega t_A=\frac{\beta ct_A}{R}. \tag{5}\]

此时 \(A\) 的位置为

\[\boldsymbol r_A=(R\sin\psi,\,R(1-\cos\psi)),\]

速度为

\[\boldsymbol v_A=\beta c(\cos\psi,\sin\psi). \tag{6}\]

设与该 \(A\) 事件在瞬时惯性系中同时的 \(B\) 事件，在 \(S\)
系中的坐标时间为 \(T_B\)。则

\[\boldsymbol r_B=(\beta cT_B,0). \tag{7}\]

两个事件在 \(A\) 的瞬时共动系中同时，意味着时间差的洛伦兹变换为零：

\[c(T_B-t_A)
-\beta\left[
\cos\psi(x_B-x_A)+\sin\psi(y_B-y_A)
\right]=0. \tag{8}\]

式（8）也可从四矢量写法得到。设两事件差为

\[\Delta x^\mu=\bigl(c(T_B-t_A),\,\boldsymbol r_B-\boldsymbol r_A\bigr),\]

\(A\) 当下的四速度为

\[u_A^\mu=\gamma\bigl(c,\,\boldsymbol v_A\bigr).\]

瞬时共动系中的同时性等价于该事件差与四速度正交：

\[u_{A\mu}\Delta x^\mu=0.\]

采用时间分量为正的度规后得到

\[c^2(T_B-t_A)-\boldsymbol v_A\cdot
(\boldsymbol r_B-\boldsymbol r_A)=0,\]

除以 \(c\)
就是式（8）。因而这是局部洛伦兹几何条件，并不是对加速观察者任意规定的一条经验公式。

先化简方括号。代入式（6）、（7）：

\[\begin{aligned}
&\cos\psi(\beta cT_B-R\sin\psi)
+\sin\psi[-R(1-\cos\psi)]\\
&=\beta cT_B\cos\psi
-R\sin\psi\cos\psi-R\sin\psi+R\sin\psi\cos\psi\\
&=\beta cT_B\cos\psi-R\sin\psi.
\end{aligned}\]

上式中两个混合项的相消可以逐项看清：

\[\cos\psi(-R\sin\psi)=-R\sin\psi\cos\psi,\]

而

\[\sin\psi\bigl[-R(1-\cos\psi)\bigr]
=-R\sin\psi+R\sin\psi\cos\psi.\]

二者相加后，\(R\sin\psi\cos\psi\) 一正一负恰好抵消，只剩位置投影
\(-R\sin\psi\)。若圆心坐标或速度方向有一个符号写错，这个简化通常就不会出现。

所以式（8）成为

\[c(T_B-t_A)-\beta
(\beta cT_B\cos\psi-R\sin\psi)=0.\]

整理所有含 \(T_B\) 的项：

\[T_B(1-\beta^2\cos\psi)
=t_A-\frac{\beta R}{c}\sin\psi. \tag{9}\]

因此

\[T_B=\frac{t_A-\dfrac{\beta R}{c}\sin\psi}
{1-\beta^2\cos\psi}. \tag{10}\]

\(B\) 的钟示数是其固有时 \(t_B=T_B/\gamma\)，故

\[\boxed{
t_B=\frac{
t_A-\dfrac{\beta R}{c}\sin\left(\dfrac{\beta ct_A}{R}\right)}
{\gamma\left[1-\beta^2\cos\left(\dfrac{\beta ct_A}{R}\right)\right]}
}. \tag{11}\]

检查 \(t_A=0\)：式（11）给出 \(t_B=0\)，与两钟在 \(O\)
点调零一致。小速度极限 \(\beta\to0\) 时，\(\gamma\to1\)，也有
\(t_B\to t_A\)。

还可作 \(\beta\) 的低阶展开。由于

\[\psi=\frac{\beta ct_A}{R}=O(\beta),\]

有

\[\sin\psi=\psi+O(\beta^3),\qquad
\cos\psi=1-\frac{\psi^2}{2}+O(\beta^4),\]

且

\[\gamma=1+\frac12\beta^2+O(\beta^4).\]

代入式（11）可见 \(t_B-t_A\) 从 \(O(\beta^2)\)
才开始出现，没有一阶修正。这与非相对论极限中同时性的绝对性相符，也为复杂分式提供了一个灵敏的符号检查。

第一问的操作顺序是：先在 \(S\) 系写出 \(A\) 的圆周世界线和 \(B\)
的直线世界线，再用洛伦兹变换确定与 \(B\) 某一事件同时的 \(A\)
事件。\(B\) 的钟示数 \(t_B\) 是固有时，对应的 \(S\) 系坐标时为
\(\gamma t_B\)。同时性条件不是令两个事件具有相同的 \(S\)
系时间，而是令它们在 \(B\)
系中的时间坐标相等，因此最终得到含三角函数的隐式方程是正常的。

第二问中先确定 \(A\) 钟示数为 \(t_A\) 时的事件 \(P_A\)。在该事件建立与
\(A\) 瞬时同速的惯性系，其同时超平面满足

\[c^2(T_B-T_A)=\boldsymbol v_A\cdot(\boldsymbol r_B-\boldsymbol r_A).\]

这里 \(\boldsymbol v_A\) 必须取事件 \(P_A\)
当下的切向速度，不能用初速度代替。把 \(B\) 世界线代入后求得的是与
\(P_A\) 同时的 \(B\) 事件在 \(S\) 系中的时间 \(T_B\)，最后还要除以
\(\gamma\) 才是 \(B\) 钟的固有时读数。

两个答案不互为简单反函数，因为``同时''依赖观察者。可用
\(t=0\)、\(\beta\to0\) 两个极限检查，并观察随后因 \(A\)
加速而出现的差别。

\newpage

\section{第 8 题}\label{ux7b2c-8-ux9898}

一个能量为 \(E_\gamma=100.0\,\mathrm{MeV}\)
的光子撞击一个初始静止、质量为 \(m_n=3.932\,\mathrm{GeV}/c^2\)
的原子核。光子被吸收后产生一对正、负电子，原子核发生反冲，反冲动量大小记为
\(p_R\)。电子静能为 \(m_ec^2=0.51\,\mathrm{MeV}\)。

\begin{enumerate}
\def\labelenumi{\arabic{enumi}.}
\tightlist
\item
  在满足能量---动量守恒的前提下，求反冲核动量 \(p_R\) 的最小值。
\item
  求反冲核速度方向与入射光子方向的最大夹角 \(\theta_M\)。
\end{enumerate}

\subsection{一、题型定位}\label{ux4e00ux9898ux578bux5b9aux4f4d-7}

把电子对看成有可变不变质量的复合体，三体末态便可按两体运动处理。用质心系动量和洛伦兹变换求反冲核的动量、角度极值，再检验阈值。

\subsection{二、知识点}\label{ux4e8cux77e5ux8bc6ux70b9-7}

\begin{itemize}
\tightlist
\item
  四动量守恒与系统不变量 \(s\)。
\item
  质心系速度与能量的洛伦兹变换。
\item
  两体末态动量公式。
\item
  速度共线变换与动量分量变换。
\item
  固定质心系动量球经洛伦兹变换后的切线极值。
\end{itemize}

\subsection{三、物理图像}\label{ux4e09ux7269ux7406ux56feux50cf-7}

实验室系的总动量沿入射光子方向。转到总质心系后，总动量为零，反冲核与电子对复合体背向运动。要使实验室系反冲核动量最小，应让质心系中的核动量尽可能大并朝质心运动反方向发射；要使偏角最大，则让洛伦兹变换后的核动量轨迹从原点看达到切线位置。

\subsection{四、解题入口}\label{ux56dbux89e3ux9898ux5165ux53e3-7}

全程以 \(\mathrm{MeV}\) 为能量单位、\(\mathrm{MeV}/c\)
为动量单位，并在中间暂取 \(c=1\)。先求初态总四动量的不变量

\[s=(m_n+E_\gamma)^2-E_\gamma^2,\]

再把电子对的最小不变质量取为
\(2m_e\)。这是反冲核在质心系中可达到最大动量的边界。

\subsection{五、解析}\label{ux4e94ux89e3ux6790-7}

图 8-1
左侧表示总质心系中电子对与反冲核动量等大反向，改变发射方向时核动量大小固定；右侧表示变回实验室系后核动量端点形成偏移曲线，最大偏角对应从原点向该曲线作切线。第（1）问先确定允许的质心系动量范围并求实验室最小值，第（2）问对固定边界动量使用切线极值条件。

\begin{figure}[H]
\centering
\includegraphics[width=0.9\linewidth,height=\textheight,keepaspectratio,alt={图 8-1　总质心系中反冲核动量方向，以及变回实验室系后的动量轨迹与最大偏角切线。}]{../figures/solution-08.pdf}
\caption*{图
8-1　总质心系中反冲核动量方向，以及变回实验室系后的动量轨迹与最大偏角切线。}
\end{figure}

下面在数值计算中暂令 \(c=1\)。原子核静能

\[m_n=3932\,\mathrm{MeV},\]

初态总能量与总动量为

\[E_t=m_n+E_\gamma=4032\,\mathrm{MeV},\]

\[p_t=E_\gamma=100.0\,\mathrm{MeV}. \tag{1}\]

总质心系相对实验室系的速度为

\[\boxed{
\beta_C=\frac{p_t}{E_t}
=\frac{100}{4032}
=2.4801587\times10^{-2}
},\]

\[\gamma_C=\frac1{\sqrt{1-\beta_C^2}}
=1.00030770. \tag{2}\]

初态不变量为

\[s=E_t^2-p_t^2,\]

也可直接从光子四动量与静止核四动量展开：

\[\begin{aligned}
s
&=(p_n+p_\gamma)^2\\
&=p_n^2+2p_n\cdot p_\gamma+p_\gamma^2\\
&=m_n^2+2m_nE_\gamma.
\end{aligned}\]

这里使用了静止核 \(p_n=(m_n,\boldsymbol0)\) 以及光子不变量
\(p_\gamma^2=0\)。数值上，

\[s=3932^2+2(3932)(100)\;\mathrm{MeV}^2,\]

与 \((4032)^2-(100)^2\)
完全相同。第二种写法通常更不容易发生两个大平方数相减造成的舍入损失。

所以质心系总能量

\[\sqrt s=\sqrt{4032^2-100^2}
=4030.75973\,\mathrm{MeV}. \tag{3}\]

\textbf{（1）反冲核动量的最小值}

按图 8-1
左侧，把电子对作为一个复合体后，末态成为两体运动；其不变质量越小，反冲核在质心系中的动量越大。变回实验室系时，最小实验室动量来自核在质心系逆着质心运动方向发射、两者几乎抵消的边界状态。

把电子对看成复合体。其不变质量最小时，正负电子在自身质心系中相对静止，故

\[M_{ee,\min}=2m_e=1.02\,\mathrm{MeV}. \tag{4}\]

此时反冲核在总质心系中的动量最大。两体末态的共同动量大小为

\[\boxed{
p_R'=\frac{
\sqrt{[s-(m_n+2m_e)^2][s-(m_n-2m_e)^2]}
}{2\sqrt s}
}. \tag{5}\]

式（5）来自两体运动学，推导如下。在总质心系中，核和复合体 \(X\)
的动量大小相同、方向相反，记为 \(p'\)：

\[E_n' + E_X'=\sqrt s,\]

\[E_n'^2=m_n^2+p'^2,\]

\[E_X'^2=M_X^2+p'^2.\]

由四动量守恒，先解核能量：

\[E_n'=\frac{s+m_n^2-M_X^2}{2\sqrt s}.\]

于是

\[\begin{aligned}
p'^2
&=E_n'^2-m_n^2\\
&=\frac{(s+m_n^2-M_X^2)^2-4sm_n^2}{4s}\\
&=\frac{[s-(m_n+M_X)^2][s-(m_n-M_X)^2]}{4s}.
\end{aligned}\]

取正平方根并令
\(M_X=2m_e\)，即得到式（5）。这里的两个因子分别体现反应阈值和两体相空间边界。

代入数据：

\[p_R'=97.54464\,\mathrm{MeV}/c. \tag{6}\]

核在质心系中的能量与速度为

\[E_n'=\sqrt{m_n^2+p_R'^2}
=3933.20975\,\mathrm{MeV},\]

\[\beta_R'=\frac{p_R'}{E_n'}
=2.4800265\times10^{-2}. \tag{7}\]

为了让实验室系核动量最小，应令核在质心系中沿 \(-x\)
方向运动，即与质心系相对实验室系的速度方向相反。共线速度变换给出实验室系核速度

\[\beta_R=\frac{\beta_C-\beta_R'}{1-\beta_C\beta_R'}. \tag{8}\]

代入式（2）、（7）：

\[\beta_R=1.3230\times10^{-6}. \tag{9}\]

这个极小量不能用低精度的 \(\beta_C-\beta_R'\) 直接估算。分子只有

\[\beta_C-\beta_R'
\approx1.32\times10^{-6},\]

而两个被减数都约为 \(2.48\times10^{-2}\)，有效数字发生明显相消。分母则为

\[1-\beta_C\beta_R'\approx0.999385,\]

虽接近
\(1\)，仍应保留到最终一步。正因为题目求的是``几乎完全抵消后的剩余量''，相对论修正和计算精度都比速度大小本身看起来更重要。

因此

\[p_{R,\min}
=\frac{m_n\beta_R}{\sqrt{1-\beta_R^2}}
\approx m_n\beta_R.\]

恢复动量单位：

\[\boxed{
p_{R,\min}\approx5.20\times10^{-3}\,\mathrm{MeV}/c
}. \tag{10}\]

该值很小，是因为 \(\beta_C\) 与反冲核可达到的 \(\beta_R'\)
极其接近，二者在反向发射时几乎抵消。

\textbf{（2）反冲方向的最大夹角}

图 8-1
右侧中，随质心系发射角变化，实验室动量端点沿一条偏移曲线运动。极角达到最大时，从动量原点指向端点的射线与轨迹相切；代数上等价于对
\(\tan\theta=p_y/p_x\) 关于发射角求极值。

设质心系中核动量与 \(+x\) 方向的夹角为 \(\vartheta'\)。其分量为

\[p_x'=p_R'\cos\vartheta',\qquad
p_y'=p_R'\sin\vartheta'.\]

变回实验室系：

\[p_x=\gamma_C(p_R'\cos\vartheta'+\beta_CE_n'),\]

\[p_y=p_R'\sin\vartheta'. \tag{11}\]

把式（11）写成几何参数形式：

\[\frac{p_x}{\gamma_C}=\beta_CE_n'+p_R'\cos\vartheta',\]

\[p_y=p_R'\sin\vartheta'.\]

消去 \(\vartheta'\) 得

\[\left(\frac{p_x}{\gamma_C}-\beta_CE_n'\right)^2+p_y^2=p_R'^2.\]

所以质心系的动量圆在实验室系中变成沿 \(x\)
方向伸缩并平移的椭圆。最大偏角是从原点向这条椭圆作切线的方向角；下面的求导条件正是这个切线几何的代数表达。

实验室系偏角满足

\[\tan\theta
=\frac{p_R'\sin\vartheta'}
{\gamma_C(p_R'\cos\vartheta'+\beta_CE_n')}. \tag{12}\]

对 \(\vartheta'\) 求极值。分母保持正时，令式（12）的导数分子为零：

\[\begin{aligned}
0={}&\cos\vartheta'
(p_R'\cos\vartheta'+\beta_CE_n')
+p_R'\sin^2\vartheta'\\
={}&p_R'+\beta_CE_n'\cos\vartheta'.
\end{aligned}\]

所以切线条件为

\[\boxed{
\cos\vartheta'=-\frac{p_R'}{\beta_CE_n'}
=-0.9999466884
}. \tag{13}\]

求导中分母的平方恒为正，真正需要令零的是

\[\begin{aligned}
N(\vartheta')
&=p_R'\cos^2\vartheta'
+\beta_CE_n'\cos\vartheta'
+p_R'\sin^2\vartheta'\\
&=p_R'+\beta_CE_n'\cos\vartheta'.
\end{aligned}\]

因此

\[\cos\vartheta'=-\frac{p_R'}{\beta_CE_n'}.\]

右端绝对值略小于 \(1\)，所以切点存在。若它大于
\(1\)，原点看不到椭圆的切线极值，最大角就必须改由参数端点判断。

将式（13）代回式（12），得到

\[\boxed{\theta_M\approx89.41^\circ}. \tag{14}\]

数值代回时建议分别保留纵、横分量：

\[p_{\perp}=p_R'\sin\vartheta',\]

\[p_{\parallel}=\gamma_C
(p_R'\cos\vartheta'+\beta_CE_n').\]

由于切点靠近 \(\vartheta'=\pi\)，\(p_{\parallel}\)
是两个近似相等量之差，而 \(p_{\perp}\) 虽小却没有同样强的相消，因此
\(p_{\perp}/p_{\parallel}\) 很大，偏角接近
\(90^\circ\)。这也解释了为什么把 \(\vartheta'\) 粗略四舍五入成
\(180^\circ\) 会完全丢失正确答案。

质心系角度为
\(\vartheta'\approx179.408^\circ\)：核几乎沿反方向发射，但保留很小的横向分量；由于实验室系纵向动量也几乎被抵消，最终形成接近
\(90^\circ\) 的大偏角。

为什么仍取电子对相对静止的边界？因为减小电子对的不变质量会增大
\(p_R'\)；式（12）所描述的实验室系动量轨迹随之向外扩张，从原点能看到的切线偏角增大，故最大偏角也在
\(M_{ee}=2m_e\) 时取得。

三体末态的自由度很多，但若把正负电子看成一个复合系统 \(X\)，反应可写成
\(\gamma+n\to n+X\)。复合系统不变质量满足
\(M_X\geqslant2m_e\)；当两电子在自身质心系中相对静止时取最小值。固定总不变量
\(s\) 后，两体质心动量由卡伦函数唯一确定，并随 \(M_X\)
减小而增大，所以两个极值都应先考察 \(M_X=2m_e\) 的边界。

最小实验室反冲动量要求质心系中的核尽量逆着质心整体运动发射。沿入射轴做洛伦兹变换后，横向动量不变，纵向动量为
\(\gamma_C(p_R'\cos\vartheta'+\beta_CE_n')\)。取 \(\vartheta'=\pi\)
给出共线候选值，还要确认它没有越过零点；若能够恰好变换到静止，则最小值应为零，本题参数下并未达到该条件。

最大夹角可把实验室末动量端点看成由质心系发射角参数化的曲线。从原点看这条曲线，方向角极值满足切线与位置矢量平行，等价于对
\(\tan\theta=p_\perp/p_\parallel\) 求导为零。得到的 \(\cos\vartheta'\)
极接近
\(-1\)，说明大角度来自纵向动量几乎抵消，而不是横向动量特别大。计算时必须统一使用能量单位，把质量写成静能
\(mc^2\) 后令 \(c=1\)，最后再恢复速度或动量单位。

\end{document}
